% =====================================================================
% Methods (body only). Can be % =====================================================================
% Methods (body only). Can be % =====================================================================
% Methods (body only). Can be % =====================================================================
% Methods (body only). Can be \input{Methods_body} in place of the
% Methods section of Manuscript_formal_20261028.tex, or compiled
% stand-alone through Methods.tex. Bibliography keys used here are
% listed in Methods_references.tex.
% Figure labels follow Manuscript_formal_20261028.tex.
% Notation follows Supplementary Note 1 (single notation system).
% =====================================================================
\section*{Methods}
\label{sec:methods}

\subsection*{Donation game, strategies and private assessment}

We consider an infinitely large, well-mixed population in which randomly
chosen donors and recipients play a donation game. A donor either cooperates,
paying a cost $c>0$ to give the recipient a benefit $b>c$, or defects, which has
no payoff consequences. Individuals use one of three strategies
$I\in\{\mathrm{ALLC},\mathrm{ALLD},\mathrm{DISC}\}$ with frequencies $f_I$,
$\sum_I f_I=1$. ALLC intends to cooperate with every recipient, ALLD intends to
defect against every recipient, and DISC intends to cooperate if and only if it
regards the recipient as good. All individuals share one social norm and one
reputation mechanism (the depths $k_1,k_2$ and readout parameters
$\alpha,\beta$ defined below); strategy evolution changes only the frequencies
$f_I$.

Reputations are private~\cite{okada2018,hilbe2018,radzvilavicius2019}: every
individual keeps its own record of every other individual. In each
interaction, an observer drawn independently of donor and recipient assesses
the donor. A DISC donor chooses its action using its own image of the
recipient, whereas the observer judges the donor's realized action using the
observer's own image of the recipient, so the two can disagree. We use
second-order norms~\cite{ohtsuki2004}. Cooperation with a recipient regarded as
good is assessed positively and defection against such a recipient negatively;
towards a recipient regarded as bad, cooperation and defection are assessed
positively with probabilities $p$ and $q$, respectively. Scoring (SC), Shunning
(SH), Stern Judging (SJ) and Simple Standing (SS) correspond to
$(p,q)=(1,0)$, $(0,0)$, $(0,1)$ and $(1,1)$.

Intended cooperation fails and becomes defection with probability $e_1$;
intended defection is always implemented. After applying the norm, the
observer records the opposite assessment with probability $e_2$. The two errors
are independent, and benefits and costs accrue only for realized cooperation.
Let $P^+_{A\mid Y}$ denote the probability that the recorded assessment is
positive, given the donor's intended action $A\in\{\mathrm C,\mathrm D\}$ and
the observer's image $Y\in\{\mathrm G,\mathrm B\}$ of the recipient. With
$\epsilon=(1-e_1)(1-e_2)+e_1e_2$,
\begin{equation}
\begin{aligned}
P^+_{\mathrm C\mid\mathrm G}&=\epsilon, &
P^+_{\mathrm C\mid\mathrm B}&=e_2+(1-2e_2)\bigl[(1-e_1)p+e_1q\bigr],\\
P^+_{\mathrm D\mid\mathrm G}&=e_2, &
P^+_{\mathrm D\mid\mathrm B}&=e_2+(1-2e_2)q
\end{aligned}
\label{eq:m-assessment}
\end{equation}
(Supplementary Note~2). Implementation error is already averaged into
Eq.~\eqref{eq:m-assessment}. If the observer regards the recipient as good with
probability $r$, intended cooperation and intended defection are therefore
assessed positively with probabilities
\begin{equation}
\Phi_{\mathrm C}(r)=rP^+_{\mathrm C\mid\mathrm G}+(1-r)P^+_{\mathrm C\mid\mathrm B},
\qquad
\Phi_{\mathrm D}(r)=rP^+_{\mathrm D\mid\mathrm G}+(1-r)P^+_{\mathrm D\mid\mathrm B},
\label{eq:m-action-assessment}
\end{equation}
and their difference, the assessment contrast, is
$\Delta\Phi(r)=\Phi_{\mathrm C}(r)-\Phi_{\mathrm D}(r)=\kappa\,[r+(p-q)(1-r)]$
with $\kappa=(1-e_1)(1-2e_2)$. Unless stated otherwise, $b=5$, $c=1$ and
$e_1=e_2=0.02$.

\subsection*{Inertial reputation records and binary readout}

Observer $i$ holds a latent record $L_{ij}$ of target $j$ in the state space
\begin{equation}
\mathcal L=\{\mathrm G,\mathrm W_1,\dots,\mathrm W_{k_1},\mathrm B,
\mathrm S_1,\dots,\mathrm S_{k_2}\},\qquad k_1,k_2\in\{0,1,2,\dots\}.
\label{eq:m-states}
\end{equation}
After each interaction, only the observer's record of the donor changes, and
only through the recorded assessment. A negative assessment moves
$\mathrm G\to\mathrm W_1\to\dots\to\mathrm W_{k_1}\to\mathrm B$ and leaves
$\mathrm B$ unchanged; a positive assessment moves
$\mathrm B\to\mathrm S_1\to\dots\to\mathrm S_{k_2}\to\mathrm G$ and leaves
$\mathrm G$ unchanged. A positive assessment returns every warning state
$\mathrm W_m$ directly to $\mathrm G$, and a negative assessment returns every
recovery state $\mathrm S_n$ directly to $\mathrm B$ (direct reset). If
$k_1=0$, a negative assessment moves $\mathrm G$ directly to $\mathrm B$; if
$k_2=0$, a positive assessment moves $\mathrm B$ directly to $\mathrm G$.
Losing good standing therefore requires $k_1+1$, and regaining it $k_2+1$,
consecutive assessments of the same target by the same observer. We call $k_1$
the tolerance depth, $k_2$ the strictness depth and $K=k_1+k_2$ the total
depth; a record has $K+2$ states (complete transition table in Supplementary Note~3 and
Supplementary Table~S5).

Actions and assessments use a binary image that is sampled from the record
whenever the record is consulted. The image is good with probability
\begin{equation}
\sigma(\mathrm G)=1,\qquad \sigma(\mathrm W_m)=\alpha^m,\qquad
\sigma(\mathrm B)=0,\qquad \sigma(\mathrm S_n)=1-\beta^n,
\qquad 0\le\alpha,\beta\le1 .
\label{eq:m-readout}
\end{equation}
Consultations are independent given the current state and do not change it.
The tolerance strength $\alpha$ sets how often a warning state is still treated
as good, and the strictness strength $\beta$ how often a recovery state is
still treated as bad. We compare three settings: no inertia,
$(k_1,k_2)=(0,0)$; full one-step inertia, $(k_1,k_2)=(1,1)$ with
$\alpha=\beta=1$; and the selected multi-step architectures with full readout,
$(2,0)$, $(3,0)$, $(4,3)$ and $(1,2)$ for SC, SH, SJ and SS (selected from
simulated reputation statistics; see below). Setting $\alpha=\beta=0$ at any
depths also yields instantaneous updating, because the image then always
equals the last recorded assessment. At one-step depths, $(\alpha,\beta)=(1,0)$
and $(0,1)$ give tolerance-only and strictness-only readout.

\subsection*{Stationary reputation map}

Consider one record whose successive recorded assessments are independent and
positive with probability $u$. Its state then follows a Markov chain with
transition matrix $T(u)=uT^{+}+(1-u)T^{-}$, where $T^{+}$ and $T^{-}$ encode the
deterministic updates after positive and negative assessments. For $0<u<1$ the
chain is irreducible and aperiodic. Write $v=1-u$. Because $\mathrm W_m$ can be
entered only from $\mathrm W_{m-1}$ (with $\mathrm W_0\equiv\mathrm G$) after a
negative assessment, and $\mathrm S_n$ only from $\mathrm S_{n-1}$ (with
$\mathrm S_0\equiv\mathrm B$) after a positive one, the stationary distribution
satisfies $\pi_{\mathrm W_m}=v^m\pi_{\mathrm G}$ and
$\pi_{\mathrm S_n}=u^n\pi_{\mathrm B}$. The only transitions from the
good-standing states $\{\mathrm G,\mathrm W_m\}$ to the bad-standing states
$\{\mathrm B,\mathrm S_n\}$ complete a run of $k_1+1$ negative assessments, and
the only reverse transitions complete a run of $k_2+1$ positive assessments.
Balancing these two fluxes gives $v^{k_1+1}\pi_{\mathrm G}=u^{k_2+1}\pi_{\mathrm B}$,
and normalization yields
\begin{equation}
\pi_{\mathrm G}=\frac{u^{k_2+1}}{Z(u)},\quad
\pi_{\mathrm W_m}=\frac{v^{m}u^{k_2+1}}{Z(u)},\quad
\pi_{\mathrm B}=\frac{v^{k_1+1}}{Z(u)},\quad
\pi_{\mathrm S_n}=\frac{u^{n}v^{k_1+1}}{Z(u)},\qquad
Z(u)=v^{k_1}+u^{k_2}-v^{k_1}u^{k_2}
\label{eq:m-stationary}
\end{equation}
(Supplementary Note~4). Averaging the readout probabilities over this
distribution defines the reputation map, the stationary probability that the
target is regarded as good when it is assessed positively with probability $u$:
\begin{equation}
\mathcal R(u)\equiv\mathcal R(u;\alpha,\beta,k_1,k_2)
=\sum_{\ell\in\mathcal L}\pi_\ell(u)\,\sigma(\ell)
=\frac{u^{k_2+1}\sum_{m=0}^{k_1}(\alpha v)^m
+v^{k_1+1}\sum_{n=1}^{k_2}u^n\bigl(1-\beta^n\bigr)}
{v^{k_1}+u^{k_2}-v^{k_1}u^{k_2}} .
\label{eq:m-reputation-map}
\end{equation}
Empty sums are zero, and $\mathcal R(0)=0$ and $\mathcal R(1)=1$ by continuity.
Without inertia, $\mathcal R(u)=u$. For one-step inertia,
$\mathcal R(u)=\{u^2[1+\alpha(1-u)]+(1-\beta)u(1-u)^2\}/(1-u+u^2)$, and with
full readout, $\mathcal R(u)=u^{k_2}(1-v^{k_1+1})/Z(u)$. The map increases with
$\alpha$ and decreases with $\beta$. Two properties are used repeatedly below:
for $k_1\ge1$, $\mathcal R'(1)=1-\alpha$, where the prime denotes $\mathrm d/\mathrm du$;
and for $k_2=0$, $\mathcal R$ is concave on $[0,1]$
(Supplementary Note~4).

\subsection*{Self-consistent private reputations}

Let $r_I$ be the probability that a randomly chosen observer regards a randomly
chosen type-$I$ individual as good, and $\bar r=\sum_I f_I r_I$. We close the
private-assessment problem with two mean-field assumptions, as in previous
analyses of private assessment~\cite{okada2018,radzvilavicius2019,kessinger2023,michelmata2024,kawakatsu2024}.
First, the donor's and the observer's images of a recipient are independent
given the recipient's strategy. Both then regard a random recipient as good
with probability $h_{\mathrm{GG}}=\sum_J f_J r_J^2$, both regard it as bad with
probability $h_{\mathrm{BB}}=\sum_J f_J(1-r_J)^2$, and each of the two ordered
disagreements occurs with probability $h_{\mathrm{GB}}=\sum_J f_J r_J(1-r_J)$.
A type-$I$ donor is therefore assessed positively with probability
\begin{equation}
\begin{aligned}
u_{\mathrm{ALLC}}&=\Phi_{\mathrm C}(\bar r),\qquad
u_{\mathrm{ALLD}}=\Phi_{\mathrm D}(\bar r),\\
u_{\mathrm{DISC}}&=h_{\mathrm{GG}}P^+_{\mathrm C\mid\mathrm G}
+h_{\mathrm{GB}}\bigl(P^+_{\mathrm C\mid\mathrm B}+P^+_{\mathrm D\mid\mathrm G}\bigr)
+h_{\mathrm{BB}}P^+_{\mathrm D\mid\mathrm B}.
\end{aligned}
\label{eq:m-strategy-assessment}
\end{equation}
Second, the assessments entering a record of a type-$I$ individual are treated
as independent with this common probability $u_I$, so that each record has
the stationary response of Eq.~\eqref{eq:m-reputation-map}. Stationary
reputations then satisfy the self-consistency equations
\begin{equation}
r_I^*=\mathcal R\bigl(u_I(\boldsymbol r^*;\boldsymbol f)\bigr),
\qquad I\in\{\mathrm{ALLC},\mathrm{ALLD},\mathrm{DISC}\}.
\label{eq:m-self-consistency}
\end{equation}
Equation~\eqref{eq:m-self-consistency} can have several solutions. Consistent
with records that start in good standing, we select at every composition the
solution reached by the auxiliary relaxation
\begin{equation}
\frac{\mathrm d r_I}{\mathrm d\tau}=\mathcal R(u_I)-r_I,\qquad r_I(0)=1 .
\label{eq:m-relaxation}
\end{equation}
The variable $\tau$ is a computational time, unrelated to evolutionary time
$t$. The fixed points of Eq.~\eqref{eq:m-relaxation} coincide with the
stationary states of the mean-field dynamics of the full distribution of
latent states, $\mathrm d\boldsymbol\pi_I/\mathrm dt'=\boldsymbol\pi_I[T(u_I)-\mathbf I]$
with $r_I=\boldsymbol\pi_I\boldsymbol\sigma$ (Supplementary Note~5). Stability
statements refer to Eq.~\eqref{eq:m-relaxation}; where checked, they agreed
with the linearization of the full latent-state system.

In a homogeneous DISC population with reputation $r$, a resident donor regards
its recipient as good with probability $r$, independently of the observer, and
Eq.~\eqref{eq:m-strategy-assessment} reduces to
\begin{equation}
u_{\mathrm R}(r)=r\,\Phi_{\mathrm C}(r)+(1-r)\,\Phi_{\mathrm D}(r).
\label{eq:m-resident-assessment}
\end{equation}
Resident reputations are the roots of
$F(r)\equiv\mathcal R[u_{\mathrm R}(r)]-r$, and their stability in
Eq.~\eqref{eq:m-relaxation} is given by the sign of $F'(r)$. A fold of a
resident branch satisfies $F(r)=F'(r)=0$, that is
$\mathcal R'[u_{\mathrm R}(r)]\,u_{\mathrm R}'(r)=1$. We denote by $r_{\mathrm R}^*$
the highest strictly stable root, which coincides with the root reached by
relaxation from $r=1$ unless the highest root is a fold, and write
$u_{\mathrm R}^*=u_{\mathrm R}(r_{\mathrm R}^*)$,
$\Phi_A^*=\Phi_A(r_{\mathrm R}^*)$ and
$\Delta\Phi^*=\Delta\Phi(r_{\mathrm R}^*)$. Resident reputations do not depend
on $b/c$.

\subsection*{Strategy evolution and cooperation}

Reputations are assumed to equilibrate quickly relative to strategy change. At
each strategy composition, payoffs are evaluated with the selected stationary
reputations. Per unit time, in which each individual acts on average once as
donor and once as recipient, the expected payoffs are
\begin{equation}
\begin{aligned}
\Pi_{\mathrm{ALLC}}&=(1-e_1)\bigl[b\,(f_{\mathrm{ALLC}}+f_{\mathrm{DISC}}r^*_{\mathrm{ALLC}})-c\bigr],\\
\Pi_{\mathrm{ALLD}}&=(1-e_1)\,b\,(f_{\mathrm{ALLC}}+f_{\mathrm{DISC}}r^*_{\mathrm{ALLD}}),\\
\Pi_{\mathrm{DISC}}&=(1-e_1)\bigl[b\,(f_{\mathrm{ALLC}}+f_{\mathrm{DISC}}r^*_{\mathrm{DISC}})-c\,\bar r^*\bigr],
\end{aligned}
\label{eq:m-payoffs}
\end{equation}
and strategy frequencies follow the replicator dynamics~\cite{hofbauer1998}
\begin{equation}
\frac{\mathrm d f_I}{\mathrm dt}=f_I\bigl(\Pi_I-\bar\Pi\bigr),\qquad
\bar\Pi=\sum_J f_J\Pi_J .
\label{eq:m-replicator}
\end{equation}
The intended and realized cooperation rates are
$\rho_{\mathrm{int}}=f_{\mathrm{ALLC}}+f_{\mathrm{DISC}}\bar r^*$ and
$\rho=(1-e_1)\rho_{\mathrm{int}}$, and $\bar\Pi=(b-c)\rho$. In a homogeneous DISC
population, $\rho_{\mathrm{int}}=r^*_{\mathrm R}$; realized cooperation
cannot exceed $1-e_1=0.98$.

To compare parameter settings globally, we integrate Eq.~\eqref{eq:m-replicator}
from $n_{\mathrm{init}}=512$ initial compositions sampled uniformly by area from
the simplex, using the same compositions for all settings that are compared.
In every scan, each trajectory ended either in the zero-cooperation set (the
ALLD vertex together with the low-reputation segment of the ALLD--DISC edge,
along which the dynamics drift slowly towards ALLD; final $\rho\le7\times10^{-4}$)
or at an attractor with positive cooperation (final $\rho\ge0.34$, and
$0.974$--$0.980$ in nearly all conditions). We therefore define the cooperative
basin as the set of initial compositions whose trajectories reach an attractor
with positive cooperation, and classify a trajectory as cooperative if its
final $\rho\ge0.01$; any threshold inside the observed gap gives the same
classification, and the gap was checked for every condition. We report the
mean long-run cooperation $\bar\rho$ (the average of the final $\rho$ over
initial compositions), the cooperative-basin size $\omega_{\mathrm C}$ (the
fraction of cooperative trajectories) and the cooperation level
$\bar\rho_{\mathrm C}$ within the basin, with
$\bar\rho\approx\omega_{\mathrm C}\bar\rho_{\mathrm C}$. For the total-depth
comparison, the envelope at total depth $K$ is the largest $\bar\rho$ among the
$K+1$ allocations with $k_1+k_2=K$.

\subsection*{Invasion of a discriminator population}

We consider a homogeneous DISC population at $r^*_{\mathrm R}$ and a rare
mutant $\mathrm X$ that uses the residents' norm and record mechanism. Because
the mutant is rare, it interacts with residents only and leaves their
reputations and assessments unchanged. Suppose that $\mathrm X$ intends to help
a resident recipient with probability $\chi_{\mathrm X}$, independently of the
observer's image of that recipient, and acts independently across interactions.
It is then assessed positively with probability
$u_{\mathrm X}=\chi_{\mathrm X}\Phi^*_{\mathrm C}+(1-\chi_{\mathrm X})\Phi^*_{\mathrm D}$,
its stationary reputation is $r^*_{\mathrm X}=\mathcal R(u_{\mathrm X})$, and its
normalized payoff advantage over residents is (Supplementary Note~7)
\begin{equation}
\Lambda_{\mathrm X}\equiv\frac{\Pi_{\mathrm X}-\Pi_{\mathrm R}}{c(1-e_1)}
=\bigl(r^*_{\mathrm R}-\chi_{\mathrm X}\bigr)
-\frac{b}{c}\bigl(r^*_{\mathrm R}-r^*_{\mathrm X}\bigr),
\label{eq:m-mutant-identity}
\end{equation}
the cooperation cost saved minus the benefit lost through reputation. The
resident itself corresponds to $\chi=r^*_{\mathrm R}$ and
$u=u^*_{\mathrm R}=r^*_{\mathrm R}\Phi^*_{\mathrm C}+(1-r^*_{\mathrm R})\Phi^*_{\mathrm D}$.

\textit{Unconditional mutants.} ALLD ($\chi=0$) has reputation
$\mathcal R(\Phi^*_{\mathrm D})$ and ALLC ($\chi=1$) has reputation
$\mathcal R(\Phi^*_{\mathrm C})$, where
$\Phi^*_{\mathrm D}=r^*_{\mathrm R}P^+_{\mathrm D\mid\mathrm G}+(1-r^*_{\mathrm R})P^+_{\mathrm D\mid\mathrm B}$
and $\Phi^*_{\mathrm C}$ analogously. When ALLD is regarded worse and ALLC
better than the residents, Eq.~\eqref{eq:m-mutant-identity} implies that the
discriminator population resists both mutants if and only if
\begin{equation}
R_{\mathrm L}\equiv
\frac{r^*_{\mathrm R}}
{r^*_{\mathrm R}-\mathcal R\bigl[r^*_{\mathrm R}P^+_{\mathrm D\mid\mathrm G}
+(1-r^*_{\mathrm R})P^+_{\mathrm D\mid\mathrm B}\bigr]}
<\frac bc<
\frac{1-r^*_{\mathrm R}}
{\mathcal R\bigl[r^*_{\mathrm R}P^+_{\mathrm C\mid\mathrm G}
+(1-r^*_{\mathrm R})P^+_{\mathrm C\mid\mathrm B}\bigr]-r^*_{\mathrm R}}
\equiv R_{\mathrm U}.
\label{eq:m-window}
\end{equation}
If ALLD is not regarded worse than the residents, no $b/c$ excludes it; if ALLC
is not regarded better, it cannot invade. Because
$u^*_{\mathrm R}=r^*_{\mathrm R}\Phi^*_{\mathrm C}+(1-r^*_{\mathrm R})\Phi^*_{\mathrm D}$,
the interval in Eq.~\eqref{eq:m-window} is non-empty if and only if
$r^*_{\mathrm R}\mathcal R(\Phi^*_{\mathrm C})+(1-r^*_{\mathrm R})\mathcal R(\Phi^*_{\mathrm D})
<\mathcal R(u^*_{\mathrm R})$, that is, if $\mathcal R$ at the resident's
assessment probability lies above the chord joining its values at
$\Phi^*_{\mathrm D}$ and $\Phi^*_{\mathrm C}$. Without inertia,
$\mathcal R$ is linear and $R_{\mathrm L}=R_{\mathrm U}=1/\Delta\Phi^*$, so no
benefit-to-cost ratio protects a discriminator population (Supplementary
Note~7). Numerically, we classify local outcomes from the signs of the
unnormalized payoff differences
$\lambda_{\mathrm{ALLD}}=\Pi_{\mathrm{ALLD}}-\Pi_{\mathrm R}$ and
$\lambda_{\mathrm{ALLC}}=\Pi_{\mathrm{ALLC}}-\Pi_{\mathrm R}$ (with $c=1$), which also
covers cases without finite thresholds: ALLD is not excluded if
$\lambda_{\mathrm{ALLD}}\ge-10^{-10}$; otherwise ALLC invades if
$\lambda_{\mathrm{ALLC}}>10^{-10}$; otherwise the residents resist both.

\textit{Opportunistic discriminators.} An opportunistic discriminator (O-DISC)
behaves like DISC, except that each intended cooperation towards a recipient it
regards as good is replaced, independently and before implementation error, by
intended defection with probability $\delta\in[0,1]$. Thus
$\chi=(1-\delta)r^*_{\mathrm R}$, and $\delta=0$ and $\delta=1$ recover DISC and
ALLD. The mutant's positive-assessment probability and reputation are
\begin{equation}
u_{\mathrm X}(\delta)=u^*_{\mathrm R}-\delta\,r^*_{\mathrm R}\,\Delta\Phi^*,
\qquad
r^*_{\mathrm X}(\delta)=\mathcal R\bigl[u_{\mathrm X}(\delta)\bigr],
\label{eq:m-odisc-reputation}
\end{equation}
and Eq.~\eqref{eq:m-mutant-identity} gives the normalized invasion payoff and
its slope at the resident,
\begin{equation}
\Lambda(\delta)=\delta\,r^*_{\mathrm R}
-\frac{b}{c}\bigl[r^*_{\mathrm R}-r^*_{\mathrm X}(\delta)\bigr],
\qquad
\Lambda'(0)=r^*_{\mathrm R}\Bigl[1-\frac{b}{c}\,\Delta\Phi^*\,
\mathcal R'(u^*_{\mathrm R})\Bigr].
\label{eq:m-odisc-invasion}
\end{equation}
A positive $\Lambda(\delta)$ means that the mutant can invade; $\Lambda(0)=0$,
and $\Lambda(1)$ is the normalized ALLD advantage. For $r^*_{\mathrm R}>0$,
$\Lambda(\delta)=\delta r^*_{\mathrm R}[1-(b/c)H(\delta)]$, where
$H(\delta)=[r^*_{\mathrm R}-r^*_{\mathrm X}(\delta)]/(\delta r^*_{\mathrm R})$ is the
reputation lost per unit of withheld cooperation, extended continuously by
$H(0)=\Delta\Phi^*\mathcal R'(u^*_{\mathrm R})$. The residents therefore exclude
every O-DISC mutant if and only if
\begin{equation}
\min_{0\le\delta\le1}\frac{r^*_{\mathrm R}-r^*_{\mathrm X}(\delta)}{\delta\,r^*_{\mathrm R}}
\;\ge\;\frac{c}{b},
\label{eq:m-security}
\end{equation}
with equality permitting neutral mutants. Because $H(\delta)$ equals
$\Delta\Phi^*$ times the slope of the chord of $\mathcal R$ between
$u_{\mathrm X}(\delta)$ and $u^*_{\mathrm R}$, for $\Delta\Phi^*>0$ the minimum
is attained at $\delta\to0$ whenever $\mathcal R$ is concave on that interval
(always for $k_2=0$), and Eq.~\eqref{eq:m-security} then reduces to
$\Lambda'(0)\le0$; in general, the whole interval must be checked. Near an
all-good resident with $k_1\ge1$, $\mathcal R'(u^*_{\mathrm R})\approx1-\alpha$,
so full tolerance readout makes weak opportunism cheap. In the error-free limit with $\beta=1$ and $k_1\ge1$, an
all-good resident excludes every O-DISC mutant if and only if
$\alpha\le1-c/b$, for every norm and every $k_2$ (Supplementary Note~8); at
$b/c=5$ this bound is the readout $\alpha=0.8$ used in Fig.~\ref{fig:opportunism}
and Extended Data Fig.~\ref{ed:opportunism-security}.

For fixed norm, depths, $\beta$ and errors, a readout $\alpha$ is secure if a
strictly stable high-reputation resident exists and satisfies
Eq.~\eqref{eq:m-security}. The lower boundary $\alpha_{\mathrm L}$ of a secure
interval, when it is not $\alpha=0$, is a fold of the high-reputation branch,
$F(r,\alpha)=\partial_rF(r,\alpha)=0$ with
$F(r,\alpha)=\mathcal R[u_{\mathrm R}(r);\alpha,\beta,k_1,k_2]-r$; the fold
itself is excluded and does not depend on $b/c$. The upper boundary
$\alpha_{\mathrm U}$, when it is not $\alpha=1$, is where
Eq.~\eqref{eq:m-security} holds with equality. In every case we computed it was
determined by $\delta\to0$, that is by
$(b/c)\,\Delta\Phi(r)\,\mathcal R'[u_{\mathrm R}(r);\alpha_{\mathrm U}]=1$
together with $F(r,\alpha_{\mathrm U})=0$. When the fold and the local condition
meet, the secure interval closes; this happens at
$b/c=u'_{\mathrm R}(r)/\Delta\Phi(r)$ evaluated at the fold, which equals 2 for
Shunning (Supplementary Note~8).

\textit{Alternative reputation filters.} For Fig.~\ref{fig:opportunism}a we
replace $\mathcal R$ by the stationary response of two alternative mechanisms
while keeping the norm, errors, resident selection and
Eqs.~\eqref{eq:m-odisc-reputation}--\eqref{eq:m-odisc-invasion} unchanged.
Aggregation, adapted from ref.~\cite{michelmata2024}, assigns a good image if at
least $Q$ of $M$ independently sampled assessments are positive,
$\mathcal R_{\mathrm{agg}}(u)=\sum_{j=Q}^{M}\binom Mj u^j(1-u)^{M-j}$, with
$(M,Q)=(2,1)$, so that $\mathcal R_{\mathrm{agg}}(u)=2u-u^2$. Quantitative
assessment, adapted from ref.~\cite{schmid2023}, keeps an integer score
$z\in\{-z_{\max},\dots,z_{\max}\}$ that moves up or down by one after a positive
or negative assessment (attempts to leave the range leave the score unchanged)
and is read as good if $z\ge z_0$; with $t=u/(1-u)$,
$\mathcal R_{\mathrm{score}}(u)=\sum_{j=z_{\max}+z_0}^{2z_{\max}}t^j\big/\sum_{j=0}^{2z_{\max}}t^j$,
and we use $z_{\max}=5$, $z_0=0$. Neither mechanism was optimized or matched in
memory to the inertial record.

\subsection*{Finite-population reputation simulations}

Individual-based simulations test reputation formation without the mean-field
closure, at fixed strategy composition; they do not simulate strategy
evolution. A population of $N$ individuals contains fixed numbers $n_I$ of each
strategy. Every DISC individual stores a latent record of every individual,
including itself, and all records start in $\mathrm G$. Records held by ALLC and
ALLD individuals are not stored because they never affect actions or
assessments. In each event: (i) a donor is drawn uniformly from the population
and a recipient uniformly from the remaining $N-1$ individuals, and an observer
is drawn independently and uniformly from the DISC individuals (the observer
may coincide with donor or recipient); (ii) the donor forms its intended
action, a DISC donor by sampling its own image of the recipient from
Eq.~\eqref{eq:m-readout}; (iii) implementation error is applied; (iv) the
observer independently samples its own image of the recipient, applies the norm
to the realized action and reverses the assessment with probability $e_2$; and
(v) the observer's record of the donor is updated. A sweep consists of
$n_{\mathrm{DISC}}N$ events, one update per stored record on average.

At sampled times we use $s_{ij}=\sigma(L_{ij})$, which averages exactly over
readout randomness. For target $j$, the mean goodness and the probability that
two distinct DISC observers give different images are
\begin{equation}
g_j=\frac{1}{n_{\mathrm{DISC}}}\sum_i s_{ij},\qquad
d_j=\binom{n_{\mathrm{DISC}}}{2}^{-1}\sum_{i<i'}
\bigl[s_{ij}(1-s_{i'j})+(1-s_{ij})s_{i'j}\bigr],
\label{eq:m-finite-estimators}
\end{equation}
where $i,i'$ run over DISC observers. Averaging over type-$I$ targets gives
$g_I$ and $d_I$, which are then averaged over sampling times. The reputational
separation and mean disagreement are
\begin{equation}
\Delta g=\frac{g_{\mathrm{DISC}}+g_{\mathrm{ALLC}}}{2}-g_{\mathrm{ALLD}},
\qquad
\bar d=\sum_I\frac{n_I}{N}\,d_I .
\label{eq:m-recognition}
\end{equation}
Unlike the mean-field theory, the simulations retain target-specific variation
and correlations between the private records of different observers.

The recognition analyses use equal strategy numbers, $\alpha=\beta=1$ and
$0\le k_1,k_2\le6$. Each run lasts 1,000 sweeps and is sampled every five sweeps
from sweep 800 to sweep 1,000 (41 samples). The depth maps in
Fig.~\ref{fig:recognition}e--h use one run per condition with $N=150$. The
replicated analysis (Extended Data Fig.~\ref{ed:recognition-selection}) uses 30
independent runs per condition with $N=150$ and reports means across runs, with
pointwise 95\% percentile intervals from 2,000 bootstrap resamples of the run
means. A condition is eligible for selection if, for both $\Delta g$ and
$\bar d$, the bootstrap interval of the difference between the second- and
first-half means of the sampling window lies within $(-0.01,0.01)$. Two slowly
relaxing SJ conditions required longer runs: $(k_1,k_2)=(4,5)$ with 3,000 sweeps
(sampled from sweep 2,400) and $(3,5)$ with 50,000 sweeps (sampled from sweep
40,000); all 196 conditions then passed. Among eligible conditions with
$\Delta g\ge0.8$ and $\bar d\le0.1$, we select the smallest $K$, breaking ties
by larger $\Delta g$, then smaller $\bar d$, $k_1$ and $k_2$. This gives the
four architectures listed above. The stars in Fig.~\ref{fig:recognition}e--h
apply the same cutoffs and tie-breaking rule to the single-run maps, without
the drift criterion, and identify the same four architectures. Selection uses
reputation statistics only.
The robustness analysis varies the separation cutoff from 0.70 to 0.90 and the
disagreement cutoff from 0.050 to 0.150 in steps of 0.05 and 0.025. The
illustrative image matrices in Fig.~\ref{fig:recognition}a--d are final states of
single runs with $N=90$ after 1,000 sweeps.

\subsection*{Numerical procedures and figure-specific settings}

All panels use $b/c=5$ and $e_1=e_2=0.02$ unless an axis varies them;
Supplementary Note~12 (Supplementary Table~S11) lists every setting, grid and
tolerance.
$\mathcal R$ and $\mathcal R'$ are evaluated in closed form. Homogeneous-DISC
roots (Fig.~\ref{fig:opportunism}; Extended Data
Figs.~\ref{ed:resident-reputation}a--d, \ref{ed:reputation-planes} and
\ref{ed:opportunism-security}) are located by evaluating $F(r)$ on a fine
uniform grid, refining sign changes by bracketing and classifying roots by the
sign of $F'(r)$. At mixed compositions, Eq.~\eqref{eq:m-relaxation} is
integrated from $r_I=1$ by fourth-order Runge--Kutta steps until
$\max_I|\mathcal R(u_I)-r_I|\le10^{-11}$ ($10^{-12}$ for
Fig.~\ref{fig:strategy-dynamics}). For local invasion (Extended Data
Figs.~\ref{ed:benefit-cost}--\ref{ed:readout-evolution}), the resident
relaxation is followed by step-limited Newton refinement, which keeps the
branch reached from $r=1$, to residual $\le10^{-13}$.

All sampled-cooperation results (Fig.~\ref{fig:memory-allocation}; Extended
Data Figs.~\ref{ed:resident-reputation}e, \ref{ed:benefit-cost},
\ref{ed:depth-evolution} and \ref{ed:readout-evolution}e--h) use one engine.
Selected reputations are tabulated exactly on a triangular lattice with 120
subdivisions and interpolated linearly only inside cells in which no reputation
component varies by more than 0.02; elsewhere, including every cell crossed by
a fold, they are computed exactly at the requested composition. Payoffs and the
replicator field are always computed from these reputations, never
interpolated. Trajectories are integrated with the explicit midpoint rule
($\Delta t=0.2$) until $\rho_{\mathrm{int}}$ changes by at most $10^{-7}$ over
each of three consecutive intervals of 100 time units, or for at most
$4\times10^4$ time units. Halving the step, refining the lattice and tightening
the tolerances changed $\bar\rho$ by at most $1.3\times10^{-5}$ and no outcome
class. Initial compositions are 512 points of a randomly shifted Halton
sequence (bases 2 and 3; seed 2026) or, for Extended Data
Fig.~\ref{ed:benefit-cost}, eight independently scrambled 512-point Sobol
sequences~\cite{virtanen2020} (seeds 2026--2033), mapped area-uniformly onto the
simplex by reflecting points of the unit square with $x+y>1$. In Extended Data
Fig.~\ref{ed:benefit-cost}, error bars are standard errors across the eight
samples and therefore describe the sampling of initial compositions; the
dynamics are deterministic. Extended Data Fig.~\ref{ed:resident-reputation}e
initializes reputations once, at $t=0$, as all-good or all-bad and then follows
them quasi-statically along each trajectory by warm-started relaxation, so that
a branch is left only where it ends.

In Fig.~\ref{fig:strategy-dynamics} and Extended Data
Fig.~\ref{ed:readout-evolution}a--d, reputations are relaxed exactly at every
point of an $84\times74$ grid clipped to the simplex. Vertex stability follows
from the two rare-strategy growth rates; edge equilibria are located on
501-point edge scans, and their markers refer to stability along the edge;
interior equilibria are classified by their Jacobian. Under SC without inertia,
all three payoffs are equal on the line
$f_{\mathrm{DISC}}=c/[b(\epsilon-e_2)]\approx0.2126$ (Supplementary Note~6).
Extended Data Fig.~\ref{ed:reputation-planes} reports the highest stable
homogeneous-DISC root and marks bistability (one-step readout plane,
$161\times161$; depth plane, $0\le k_1,k_2\le8$); for every root in the depth
planes, stability in Eq.~\eqref{eq:m-relaxation} agreed with the linearization of
the full latent-state dynamics. On the ALLD--DISC edge (Extended Data
Fig.~\ref{ed:edge-geometry}), Eq.~\eqref{eq:m-self-consistency} reduces exactly
to one equation in $\bar r$, which is solved for 1,071 values of
$f_{\mathrm{DISC}}\in[0.01,1]$. Local invasion maps use $241\times241$ points of
$e_1,e_2\in[0,0.45]$ (Extended Data Fig.~\ref{ed:error-invasion}) and
$801\times801$ points of $(\alpha,\beta)$ (Extended Data
Fig.~\ref{ed:readout-evolution}i--l). In Fig.~\ref{fig:opportunism},
$\Lambda(\delta)$ is evaluated at 5,015 values of $\delta$, concentrated near
$\delta=0$, with sign changes refined by Brent's method; $\delta_{\mathrm c}$ is
the upper end of the interval with $\Lambda>0$. Residents of the alternative
filters are obtained from $r=1$ by damped iteration,
$r\leftarrow0.8r+0.2\,\mathcal R_{\mathrm X}[u_{\mathrm R}(r)]$ with
$\mathcal R_{\mathrm X}\in\{\mathcal R_{\mathrm{agg}},\mathcal R_{\mathrm{score}}\}$.
For Extended Data Fig.~\ref{ed:opportunism-security} ($\ge601$ values of
$\alpha$, 161 values of $b/c\in[2,8]$, $\beta=1$), $\min_\delta H(\delta)$ is
found by comparing $H$ at $\delta=0$, $\delta=1$ and all real roots in
$(10^{-8},1)$ of the numerator of the rational function
$\mathrm dH/\mathrm d\delta$, and a readout is classified as secure if
$1-(b/c)\min_\delta H\le10^{-10}$. Fold endpoints are excluded; endpoints at the
domain boundary or at neutral invasion are included.

\subsection*{Data and code availability}

The code, numerical data and figure scripts used in this study are provided
with the paper (verified computational package; one notebook per figure, a C
engine for the strategy dynamics, input configurations, validation records and
checksums). \textit{[Repository and DOI to be added before submission.]}
 in place of the
% Methods section of Manuscript_formal_20261028.tex, or compiled
% stand-alone through Methods.tex. Bibliography keys used here are
% listed in Methods_references.tex.
% Figure labels follow Manuscript_formal_20261028.tex.
% Notation follows Supplementary Note 1 (single notation system).
% =====================================================================
\section*{Methods}
\label{sec:methods}

\subsection*{Donation game, strategies and private assessment}

We consider an infinitely large, well-mixed population in which randomly
chosen donors and recipients play a donation game. A donor either cooperates,
paying a cost $c>0$ to give the recipient a benefit $b>c$, or defects, which has
no payoff consequences. Individuals use one of three strategies
$I\in\{\mathrm{ALLC},\mathrm{ALLD},\mathrm{DISC}\}$ with frequencies $f_I$,
$\sum_I f_I=1$. ALLC intends to cooperate with every recipient, ALLD intends to
defect against every recipient, and DISC intends to cooperate if and only if it
regards the recipient as good. All individuals share one social norm and one
reputation mechanism (the depths $k_1,k_2$ and readout parameters
$\alpha,\beta$ defined below); strategy evolution changes only the frequencies
$f_I$.

Reputations are private~\cite{okada2018,hilbe2018,radzvilavicius2019}: every
individual keeps its own record of every other individual. In each
interaction, an observer drawn independently of donor and recipient assesses
the donor. A DISC donor chooses its action using its own image of the
recipient, whereas the observer judges the donor's realized action using the
observer's own image of the recipient, so the two can disagree. We use
second-order norms~\cite{ohtsuki2004}. Cooperation with a recipient regarded as
good is assessed positively and defection against such a recipient negatively;
towards a recipient regarded as bad, cooperation and defection are assessed
positively with probabilities $p$ and $q$, respectively. Scoring (SC), Shunning
(SH), Stern Judging (SJ) and Simple Standing (SS) correspond to
$(p,q)=(1,0)$, $(0,0)$, $(0,1)$ and $(1,1)$.

Intended cooperation fails and becomes defection with probability $e_1$;
intended defection is always implemented. After applying the norm, the
observer records the opposite assessment with probability $e_2$. The two errors
are independent, and benefits and costs accrue only for realized cooperation.
Let $P^+_{A\mid Y}$ denote the probability that the recorded assessment is
positive, given the donor's intended action $A\in\{\mathrm C,\mathrm D\}$ and
the observer's image $Y\in\{\mathrm G,\mathrm B\}$ of the recipient. With
$\epsilon=(1-e_1)(1-e_2)+e_1e_2$,
\begin{equation}
\begin{aligned}
P^+_{\mathrm C\mid\mathrm G}&=\epsilon, &
P^+_{\mathrm C\mid\mathrm B}&=e_2+(1-2e_2)\bigl[(1-e_1)p+e_1q\bigr],\\
P^+_{\mathrm D\mid\mathrm G}&=e_2, &
P^+_{\mathrm D\mid\mathrm B}&=e_2+(1-2e_2)q
\end{aligned}
\label{eq:m-assessment}
\end{equation}
(Supplementary Note~2). Implementation error is already averaged into
Eq.~\eqref{eq:m-assessment}. If the observer regards the recipient as good with
probability $r$, intended cooperation and intended defection are therefore
assessed positively with probabilities
\begin{equation}
\Phi_{\mathrm C}(r)=rP^+_{\mathrm C\mid\mathrm G}+(1-r)P^+_{\mathrm C\mid\mathrm B},
\qquad
\Phi_{\mathrm D}(r)=rP^+_{\mathrm D\mid\mathrm G}+(1-r)P^+_{\mathrm D\mid\mathrm B},
\label{eq:m-action-assessment}
\end{equation}
and their difference, the assessment contrast, is
$\Delta\Phi(r)=\Phi_{\mathrm C}(r)-\Phi_{\mathrm D}(r)=\kappa\,[r+(p-q)(1-r)]$
with $\kappa=(1-e_1)(1-2e_2)$. Unless stated otherwise, $b=5$, $c=1$ and
$e_1=e_2=0.02$.

\subsection*{Inertial reputation records and binary readout}

Observer $i$ holds a latent record $L_{ij}$ of target $j$ in the state space
\begin{equation}
\mathcal L=\{\mathrm G,\mathrm W_1,\dots,\mathrm W_{k_1},\mathrm B,
\mathrm S_1,\dots,\mathrm S_{k_2}\},\qquad k_1,k_2\in\{0,1,2,\dots\}.
\label{eq:m-states}
\end{equation}
After each interaction, only the observer's record of the donor changes, and
only through the recorded assessment. A negative assessment moves
$\mathrm G\to\mathrm W_1\to\dots\to\mathrm W_{k_1}\to\mathrm B$ and leaves
$\mathrm B$ unchanged; a positive assessment moves
$\mathrm B\to\mathrm S_1\to\dots\to\mathrm S_{k_2}\to\mathrm G$ and leaves
$\mathrm G$ unchanged. A positive assessment returns every warning state
$\mathrm W_m$ directly to $\mathrm G$, and a negative assessment returns every
recovery state $\mathrm S_n$ directly to $\mathrm B$ (direct reset). If
$k_1=0$, a negative assessment moves $\mathrm G$ directly to $\mathrm B$; if
$k_2=0$, a positive assessment moves $\mathrm B$ directly to $\mathrm G$.
Losing good standing therefore requires $k_1+1$, and regaining it $k_2+1$,
consecutive assessments of the same target by the same observer. We call $k_1$
the tolerance depth, $k_2$ the strictness depth and $K=k_1+k_2$ the total
depth; a record has $K+2$ states (complete transition table in Supplementary Note~3 and
Supplementary Table~S5).

Actions and assessments use a binary image that is sampled from the record
whenever the record is consulted. The image is good with probability
\begin{equation}
\sigma(\mathrm G)=1,\qquad \sigma(\mathrm W_m)=\alpha^m,\qquad
\sigma(\mathrm B)=0,\qquad \sigma(\mathrm S_n)=1-\beta^n,
\qquad 0\le\alpha,\beta\le1 .
\label{eq:m-readout}
\end{equation}
Consultations are independent given the current state and do not change it.
The tolerance strength $\alpha$ sets how often a warning state is still treated
as good, and the strictness strength $\beta$ how often a recovery state is
still treated as bad. We compare three settings: no inertia,
$(k_1,k_2)=(0,0)$; full one-step inertia, $(k_1,k_2)=(1,1)$ with
$\alpha=\beta=1$; and the selected multi-step architectures with full readout,
$(2,0)$, $(3,0)$, $(4,3)$ and $(1,2)$ for SC, SH, SJ and SS (selected from
simulated reputation statistics; see below). Setting $\alpha=\beta=0$ at any
depths also yields instantaneous updating, because the image then always
equals the last recorded assessment. At one-step depths, $(\alpha,\beta)=(1,0)$
and $(0,1)$ give tolerance-only and strictness-only readout.

\subsection*{Stationary reputation map}

Consider one record whose successive recorded assessments are independent and
positive with probability $u$. Its state then follows a Markov chain with
transition matrix $T(u)=uT^{+}+(1-u)T^{-}$, where $T^{+}$ and $T^{-}$ encode the
deterministic updates after positive and negative assessments. For $0<u<1$ the
chain is irreducible and aperiodic. Write $v=1-u$. Because $\mathrm W_m$ can be
entered only from $\mathrm W_{m-1}$ (with $\mathrm W_0\equiv\mathrm G$) after a
negative assessment, and $\mathrm S_n$ only from $\mathrm S_{n-1}$ (with
$\mathrm S_0\equiv\mathrm B$) after a positive one, the stationary distribution
satisfies $\pi_{\mathrm W_m}=v^m\pi_{\mathrm G}$ and
$\pi_{\mathrm S_n}=u^n\pi_{\mathrm B}$. The only transitions from the
good-standing states $\{\mathrm G,\mathrm W_m\}$ to the bad-standing states
$\{\mathrm B,\mathrm S_n\}$ complete a run of $k_1+1$ negative assessments, and
the only reverse transitions complete a run of $k_2+1$ positive assessments.
Balancing these two fluxes gives $v^{k_1+1}\pi_{\mathrm G}=u^{k_2+1}\pi_{\mathrm B}$,
and normalization yields
\begin{equation}
\pi_{\mathrm G}=\frac{u^{k_2+1}}{Z(u)},\quad
\pi_{\mathrm W_m}=\frac{v^{m}u^{k_2+1}}{Z(u)},\quad
\pi_{\mathrm B}=\frac{v^{k_1+1}}{Z(u)},\quad
\pi_{\mathrm S_n}=\frac{u^{n}v^{k_1+1}}{Z(u)},\qquad
Z(u)=v^{k_1}+u^{k_2}-v^{k_1}u^{k_2}
\label{eq:m-stationary}
\end{equation}
(Supplementary Note~4). Averaging the readout probabilities over this
distribution defines the reputation map, the stationary probability that the
target is regarded as good when it is assessed positively with probability $u$:
\begin{equation}
\mathcal R(u)\equiv\mathcal R(u;\alpha,\beta,k_1,k_2)
=\sum_{\ell\in\mathcal L}\pi_\ell(u)\,\sigma(\ell)
=\frac{u^{k_2+1}\sum_{m=0}^{k_1}(\alpha v)^m
+v^{k_1+1}\sum_{n=1}^{k_2}u^n\bigl(1-\beta^n\bigr)}
{v^{k_1}+u^{k_2}-v^{k_1}u^{k_2}} .
\label{eq:m-reputation-map}
\end{equation}
Empty sums are zero, and $\mathcal R(0)=0$ and $\mathcal R(1)=1$ by continuity.
Without inertia, $\mathcal R(u)=u$. For one-step inertia,
$\mathcal R(u)=\{u^2[1+\alpha(1-u)]+(1-\beta)u(1-u)^2\}/(1-u+u^2)$, and with
full readout, $\mathcal R(u)=u^{k_2}(1-v^{k_1+1})/Z(u)$. The map increases with
$\alpha$ and decreases with $\beta$. Two properties are used repeatedly below:
for $k_1\ge1$, $\mathcal R'(1)=1-\alpha$, where the prime denotes $\mathrm d/\mathrm du$;
and for $k_2=0$, $\mathcal R$ is concave on $[0,1]$
(Supplementary Note~4).

\subsection*{Self-consistent private reputations}

Let $r_I$ be the probability that a randomly chosen observer regards a randomly
chosen type-$I$ individual as good, and $\bar r=\sum_I f_I r_I$. We close the
private-assessment problem with two mean-field assumptions, as in previous
analyses of private assessment~\cite{okada2018,radzvilavicius2019,kessinger2023,michelmata2024,kawakatsu2024}.
First, the donor's and the observer's images of a recipient are independent
given the recipient's strategy. Both then regard a random recipient as good
with probability $h_{\mathrm{GG}}=\sum_J f_J r_J^2$, both regard it as bad with
probability $h_{\mathrm{BB}}=\sum_J f_J(1-r_J)^2$, and each of the two ordered
disagreements occurs with probability $h_{\mathrm{GB}}=\sum_J f_J r_J(1-r_J)$.
A type-$I$ donor is therefore assessed positively with probability
\begin{equation}
\begin{aligned}
u_{\mathrm{ALLC}}&=\Phi_{\mathrm C}(\bar r),\qquad
u_{\mathrm{ALLD}}=\Phi_{\mathrm D}(\bar r),\\
u_{\mathrm{DISC}}&=h_{\mathrm{GG}}P^+_{\mathrm C\mid\mathrm G}
+h_{\mathrm{GB}}\bigl(P^+_{\mathrm C\mid\mathrm B}+P^+_{\mathrm D\mid\mathrm G}\bigr)
+h_{\mathrm{BB}}P^+_{\mathrm D\mid\mathrm B}.
\end{aligned}
\label{eq:m-strategy-assessment}
\end{equation}
Second, the assessments entering a record of a type-$I$ individual are treated
as independent with this common probability $u_I$, so that each record has
the stationary response of Eq.~\eqref{eq:m-reputation-map}. Stationary
reputations then satisfy the self-consistency equations
\begin{equation}
r_I^*=\mathcal R\bigl(u_I(\boldsymbol r^*;\boldsymbol f)\bigr),
\qquad I\in\{\mathrm{ALLC},\mathrm{ALLD},\mathrm{DISC}\}.
\label{eq:m-self-consistency}
\end{equation}
Equation~\eqref{eq:m-self-consistency} can have several solutions. Consistent
with records that start in good standing, we select at every composition the
solution reached by the auxiliary relaxation
\begin{equation}
\frac{\mathrm d r_I}{\mathrm d\tau}=\mathcal R(u_I)-r_I,\qquad r_I(0)=1 .
\label{eq:m-relaxation}
\end{equation}
The variable $\tau$ is a computational time, unrelated to evolutionary time
$t$. The fixed points of Eq.~\eqref{eq:m-relaxation} coincide with the
stationary states of the mean-field dynamics of the full distribution of
latent states, $\mathrm d\boldsymbol\pi_I/\mathrm dt'=\boldsymbol\pi_I[T(u_I)-\mathbf I]$
with $r_I=\boldsymbol\pi_I\boldsymbol\sigma$ (Supplementary Note~5). Stability
statements refer to Eq.~\eqref{eq:m-relaxation}; where checked, they agreed
with the linearization of the full latent-state system.

In a homogeneous DISC population with reputation $r$, a resident donor regards
its recipient as good with probability $r$, independently of the observer, and
Eq.~\eqref{eq:m-strategy-assessment} reduces to
\begin{equation}
u_{\mathrm R}(r)=r\,\Phi_{\mathrm C}(r)+(1-r)\,\Phi_{\mathrm D}(r).
\label{eq:m-resident-assessment}
\end{equation}
Resident reputations are the roots of
$F(r)\equiv\mathcal R[u_{\mathrm R}(r)]-r$, and their stability in
Eq.~\eqref{eq:m-relaxation} is given by the sign of $F'(r)$. A fold of a
resident branch satisfies $F(r)=F'(r)=0$, that is
$\mathcal R'[u_{\mathrm R}(r)]\,u_{\mathrm R}'(r)=1$. We denote by $r_{\mathrm R}^*$
the highest strictly stable root, which coincides with the root reached by
relaxation from $r=1$ unless the highest root is a fold, and write
$u_{\mathrm R}^*=u_{\mathrm R}(r_{\mathrm R}^*)$,
$\Phi_A^*=\Phi_A(r_{\mathrm R}^*)$ and
$\Delta\Phi^*=\Delta\Phi(r_{\mathrm R}^*)$. Resident reputations do not depend
on $b/c$.

\subsection*{Strategy evolution and cooperation}

Reputations are assumed to equilibrate quickly relative to strategy change. At
each strategy composition, payoffs are evaluated with the selected stationary
reputations. Per unit time, in which each individual acts on average once as
donor and once as recipient, the expected payoffs are
\begin{equation}
\begin{aligned}
\Pi_{\mathrm{ALLC}}&=(1-e_1)\bigl[b\,(f_{\mathrm{ALLC}}+f_{\mathrm{DISC}}r^*_{\mathrm{ALLC}})-c\bigr],\\
\Pi_{\mathrm{ALLD}}&=(1-e_1)\,b\,(f_{\mathrm{ALLC}}+f_{\mathrm{DISC}}r^*_{\mathrm{ALLD}}),\\
\Pi_{\mathrm{DISC}}&=(1-e_1)\bigl[b\,(f_{\mathrm{ALLC}}+f_{\mathrm{DISC}}r^*_{\mathrm{DISC}})-c\,\bar r^*\bigr],
\end{aligned}
\label{eq:m-payoffs}
\end{equation}
and strategy frequencies follow the replicator dynamics~\cite{hofbauer1998}
\begin{equation}
\frac{\mathrm d f_I}{\mathrm dt}=f_I\bigl(\Pi_I-\bar\Pi\bigr),\qquad
\bar\Pi=\sum_J f_J\Pi_J .
\label{eq:m-replicator}
\end{equation}
The intended and realized cooperation rates are
$\rho_{\mathrm{int}}=f_{\mathrm{ALLC}}+f_{\mathrm{DISC}}\bar r^*$ and
$\rho=(1-e_1)\rho_{\mathrm{int}}$, and $\bar\Pi=(b-c)\rho$. In a homogeneous DISC
population, $\rho_{\mathrm{int}}=r^*_{\mathrm R}$; realized cooperation
cannot exceed $1-e_1=0.98$.

To compare parameter settings globally, we integrate Eq.~\eqref{eq:m-replicator}
from $n_{\mathrm{init}}=512$ initial compositions sampled uniformly by area from
the simplex, using the same compositions for all settings that are compared.
In every scan, each trajectory ended either in the zero-cooperation set (the
ALLD vertex together with the low-reputation segment of the ALLD--DISC edge,
along which the dynamics drift slowly towards ALLD; final $\rho\le7\times10^{-4}$)
or at an attractor with positive cooperation (final $\rho\ge0.34$, and
$0.974$--$0.980$ in nearly all conditions). We therefore define the cooperative
basin as the set of initial compositions whose trajectories reach an attractor
with positive cooperation, and classify a trajectory as cooperative if its
final $\rho\ge0.01$; any threshold inside the observed gap gives the same
classification, and the gap was checked for every condition. We report the
mean long-run cooperation $\bar\rho$ (the average of the final $\rho$ over
initial compositions), the cooperative-basin size $\omega_{\mathrm C}$ (the
fraction of cooperative trajectories) and the cooperation level
$\bar\rho_{\mathrm C}$ within the basin, with
$\bar\rho\approx\omega_{\mathrm C}\bar\rho_{\mathrm C}$. For the total-depth
comparison, the envelope at total depth $K$ is the largest $\bar\rho$ among the
$K+1$ allocations with $k_1+k_2=K$.

\subsection*{Invasion of a discriminator population}

We consider a homogeneous DISC population at $r^*_{\mathrm R}$ and a rare
mutant $\mathrm X$ that uses the residents' norm and record mechanism. Because
the mutant is rare, it interacts with residents only and leaves their
reputations and assessments unchanged. Suppose that $\mathrm X$ intends to help
a resident recipient with probability $\chi_{\mathrm X}$, independently of the
observer's image of that recipient, and acts independently across interactions.
It is then assessed positively with probability
$u_{\mathrm X}=\chi_{\mathrm X}\Phi^*_{\mathrm C}+(1-\chi_{\mathrm X})\Phi^*_{\mathrm D}$,
its stationary reputation is $r^*_{\mathrm X}=\mathcal R(u_{\mathrm X})$, and its
normalized payoff advantage over residents is (Supplementary Note~7)
\begin{equation}
\Lambda_{\mathrm X}\equiv\frac{\Pi_{\mathrm X}-\Pi_{\mathrm R}}{c(1-e_1)}
=\bigl(r^*_{\mathrm R}-\chi_{\mathrm X}\bigr)
-\frac{b}{c}\bigl(r^*_{\mathrm R}-r^*_{\mathrm X}\bigr),
\label{eq:m-mutant-identity}
\end{equation}
the cooperation cost saved minus the benefit lost through reputation. The
resident itself corresponds to $\chi=r^*_{\mathrm R}$ and
$u=u^*_{\mathrm R}=r^*_{\mathrm R}\Phi^*_{\mathrm C}+(1-r^*_{\mathrm R})\Phi^*_{\mathrm D}$.

\textit{Unconditional mutants.} ALLD ($\chi=0$) has reputation
$\mathcal R(\Phi^*_{\mathrm D})$ and ALLC ($\chi=1$) has reputation
$\mathcal R(\Phi^*_{\mathrm C})$, where
$\Phi^*_{\mathrm D}=r^*_{\mathrm R}P^+_{\mathrm D\mid\mathrm G}+(1-r^*_{\mathrm R})P^+_{\mathrm D\mid\mathrm B}$
and $\Phi^*_{\mathrm C}$ analogously. When ALLD is regarded worse and ALLC
better than the residents, Eq.~\eqref{eq:m-mutant-identity} implies that the
discriminator population resists both mutants if and only if
\begin{equation}
R_{\mathrm L}\equiv
\frac{r^*_{\mathrm R}}
{r^*_{\mathrm R}-\mathcal R\bigl[r^*_{\mathrm R}P^+_{\mathrm D\mid\mathrm G}
+(1-r^*_{\mathrm R})P^+_{\mathrm D\mid\mathrm B}\bigr]}
<\frac bc<
\frac{1-r^*_{\mathrm R}}
{\mathcal R\bigl[r^*_{\mathrm R}P^+_{\mathrm C\mid\mathrm G}
+(1-r^*_{\mathrm R})P^+_{\mathrm C\mid\mathrm B}\bigr]-r^*_{\mathrm R}}
\equiv R_{\mathrm U}.
\label{eq:m-window}
\end{equation}
If ALLD is not regarded worse than the residents, no $b/c$ excludes it; if ALLC
is not regarded better, it cannot invade. Because
$u^*_{\mathrm R}=r^*_{\mathrm R}\Phi^*_{\mathrm C}+(1-r^*_{\mathrm R})\Phi^*_{\mathrm D}$,
the interval in Eq.~\eqref{eq:m-window} is non-empty if and only if
$r^*_{\mathrm R}\mathcal R(\Phi^*_{\mathrm C})+(1-r^*_{\mathrm R})\mathcal R(\Phi^*_{\mathrm D})
<\mathcal R(u^*_{\mathrm R})$, that is, if $\mathcal R$ at the resident's
assessment probability lies above the chord joining its values at
$\Phi^*_{\mathrm D}$ and $\Phi^*_{\mathrm C}$. Without inertia,
$\mathcal R$ is linear and $R_{\mathrm L}=R_{\mathrm U}=1/\Delta\Phi^*$, so no
benefit-to-cost ratio protects a discriminator population (Supplementary
Note~7). Numerically, we classify local outcomes from the signs of the
unnormalized payoff differences
$\lambda_{\mathrm{ALLD}}=\Pi_{\mathrm{ALLD}}-\Pi_{\mathrm R}$ and
$\lambda_{\mathrm{ALLC}}=\Pi_{\mathrm{ALLC}}-\Pi_{\mathrm R}$ (with $c=1$), which also
covers cases without finite thresholds: ALLD is not excluded if
$\lambda_{\mathrm{ALLD}}\ge-10^{-10}$; otherwise ALLC invades if
$\lambda_{\mathrm{ALLC}}>10^{-10}$; otherwise the residents resist both.

\textit{Opportunistic discriminators.} An opportunistic discriminator (O-DISC)
behaves like DISC, except that each intended cooperation towards a recipient it
regards as good is replaced, independently and before implementation error, by
intended defection with probability $\delta\in[0,1]$. Thus
$\chi=(1-\delta)r^*_{\mathrm R}$, and $\delta=0$ and $\delta=1$ recover DISC and
ALLD. The mutant's positive-assessment probability and reputation are
\begin{equation}
u_{\mathrm X}(\delta)=u^*_{\mathrm R}-\delta\,r^*_{\mathrm R}\,\Delta\Phi^*,
\qquad
r^*_{\mathrm X}(\delta)=\mathcal R\bigl[u_{\mathrm X}(\delta)\bigr],
\label{eq:m-odisc-reputation}
\end{equation}
and Eq.~\eqref{eq:m-mutant-identity} gives the normalized invasion payoff and
its slope at the resident,
\begin{equation}
\Lambda(\delta)=\delta\,r^*_{\mathrm R}
-\frac{b}{c}\bigl[r^*_{\mathrm R}-r^*_{\mathrm X}(\delta)\bigr],
\qquad
\Lambda'(0)=r^*_{\mathrm R}\Bigl[1-\frac{b}{c}\,\Delta\Phi^*\,
\mathcal R'(u^*_{\mathrm R})\Bigr].
\label{eq:m-odisc-invasion}
\end{equation}
A positive $\Lambda(\delta)$ means that the mutant can invade; $\Lambda(0)=0$,
and $\Lambda(1)$ is the normalized ALLD advantage. For $r^*_{\mathrm R}>0$,
$\Lambda(\delta)=\delta r^*_{\mathrm R}[1-(b/c)H(\delta)]$, where
$H(\delta)=[r^*_{\mathrm R}-r^*_{\mathrm X}(\delta)]/(\delta r^*_{\mathrm R})$ is the
reputation lost per unit of withheld cooperation, extended continuously by
$H(0)=\Delta\Phi^*\mathcal R'(u^*_{\mathrm R})$. The residents therefore exclude
every O-DISC mutant if and only if
\begin{equation}
\min_{0\le\delta\le1}\frac{r^*_{\mathrm R}-r^*_{\mathrm X}(\delta)}{\delta\,r^*_{\mathrm R}}
\;\ge\;\frac{c}{b},
\label{eq:m-security}
\end{equation}
with equality permitting neutral mutants. Because $H(\delta)$ equals
$\Delta\Phi^*$ times the slope of the chord of $\mathcal R$ between
$u_{\mathrm X}(\delta)$ and $u^*_{\mathrm R}$, for $\Delta\Phi^*>0$ the minimum
is attained at $\delta\to0$ whenever $\mathcal R$ is concave on that interval
(always for $k_2=0$), and Eq.~\eqref{eq:m-security} then reduces to
$\Lambda'(0)\le0$; in general, the whole interval must be checked. Near an
all-good resident with $k_1\ge1$, $\mathcal R'(u^*_{\mathrm R})\approx1-\alpha$,
so full tolerance readout makes weak opportunism cheap. In the error-free limit with $\beta=1$ and $k_1\ge1$, an
all-good resident excludes every O-DISC mutant if and only if
$\alpha\le1-c/b$, for every norm and every $k_2$ (Supplementary Note~8); at
$b/c=5$ this bound is the readout $\alpha=0.8$ used in Fig.~\ref{fig:opportunism}
and Extended Data Fig.~\ref{ed:opportunism-security}.

For fixed norm, depths, $\beta$ and errors, a readout $\alpha$ is secure if a
strictly stable high-reputation resident exists and satisfies
Eq.~\eqref{eq:m-security}. The lower boundary $\alpha_{\mathrm L}$ of a secure
interval, when it is not $\alpha=0$, is a fold of the high-reputation branch,
$F(r,\alpha)=\partial_rF(r,\alpha)=0$ with
$F(r,\alpha)=\mathcal R[u_{\mathrm R}(r);\alpha,\beta,k_1,k_2]-r$; the fold
itself is excluded and does not depend on $b/c$. The upper boundary
$\alpha_{\mathrm U}$, when it is not $\alpha=1$, is where
Eq.~\eqref{eq:m-security} holds with equality. In every case we computed it was
determined by $\delta\to0$, that is by
$(b/c)\,\Delta\Phi(r)\,\mathcal R'[u_{\mathrm R}(r);\alpha_{\mathrm U}]=1$
together with $F(r,\alpha_{\mathrm U})=0$. When the fold and the local condition
meet, the secure interval closes; this happens at
$b/c=u'_{\mathrm R}(r)/\Delta\Phi(r)$ evaluated at the fold, which equals 2 for
Shunning (Supplementary Note~8).

\textit{Alternative reputation filters.} For Fig.~\ref{fig:opportunism}a we
replace $\mathcal R$ by the stationary response of two alternative mechanisms
while keeping the norm, errors, resident selection and
Eqs.~\eqref{eq:m-odisc-reputation}--\eqref{eq:m-odisc-invasion} unchanged.
Aggregation, adapted from ref.~\cite{michelmata2024}, assigns a good image if at
least $Q$ of $M$ independently sampled assessments are positive,
$\mathcal R_{\mathrm{agg}}(u)=\sum_{j=Q}^{M}\binom Mj u^j(1-u)^{M-j}$, with
$(M,Q)=(2,1)$, so that $\mathcal R_{\mathrm{agg}}(u)=2u-u^2$. Quantitative
assessment, adapted from ref.~\cite{schmid2023}, keeps an integer score
$z\in\{-z_{\max},\dots,z_{\max}\}$ that moves up or down by one after a positive
or negative assessment (attempts to leave the range leave the score unchanged)
and is read as good if $z\ge z_0$; with $t=u/(1-u)$,
$\mathcal R_{\mathrm{score}}(u)=\sum_{j=z_{\max}+z_0}^{2z_{\max}}t^j\big/\sum_{j=0}^{2z_{\max}}t^j$,
and we use $z_{\max}=5$, $z_0=0$. Neither mechanism was optimized or matched in
memory to the inertial record.

\subsection*{Finite-population reputation simulations}

Individual-based simulations test reputation formation without the mean-field
closure, at fixed strategy composition; they do not simulate strategy
evolution. A population of $N$ individuals contains fixed numbers $n_I$ of each
strategy. Every DISC individual stores a latent record of every individual,
including itself, and all records start in $\mathrm G$. Records held by ALLC and
ALLD individuals are not stored because they never affect actions or
assessments. In each event: (i) a donor is drawn uniformly from the population
and a recipient uniformly from the remaining $N-1$ individuals, and an observer
is drawn independently and uniformly from the DISC individuals (the observer
may coincide with donor or recipient); (ii) the donor forms its intended
action, a DISC donor by sampling its own image of the recipient from
Eq.~\eqref{eq:m-readout}; (iii) implementation error is applied; (iv) the
observer independently samples its own image of the recipient, applies the norm
to the realized action and reverses the assessment with probability $e_2$; and
(v) the observer's record of the donor is updated. A sweep consists of
$n_{\mathrm{DISC}}N$ events, one update per stored record on average.

At sampled times we use $s_{ij}=\sigma(L_{ij})$, which averages exactly over
readout randomness. For target $j$, the mean goodness and the probability that
two distinct DISC observers give different images are
\begin{equation}
g_j=\frac{1}{n_{\mathrm{DISC}}}\sum_i s_{ij},\qquad
d_j=\binom{n_{\mathrm{DISC}}}{2}^{-1}\sum_{i<i'}
\bigl[s_{ij}(1-s_{i'j})+(1-s_{ij})s_{i'j}\bigr],
\label{eq:m-finite-estimators}
\end{equation}
where $i,i'$ run over DISC observers. Averaging over type-$I$ targets gives
$g_I$ and $d_I$, which are then averaged over sampling times. The reputational
separation and mean disagreement are
\begin{equation}
\Delta g=\frac{g_{\mathrm{DISC}}+g_{\mathrm{ALLC}}}{2}-g_{\mathrm{ALLD}},
\qquad
\bar d=\sum_I\frac{n_I}{N}\,d_I .
\label{eq:m-recognition}
\end{equation}
Unlike the mean-field theory, the simulations retain target-specific variation
and correlations between the private records of different observers.

The recognition analyses use equal strategy numbers, $\alpha=\beta=1$ and
$0\le k_1,k_2\le6$. Each run lasts 1,000 sweeps and is sampled every five sweeps
from sweep 800 to sweep 1,000 (41 samples). The depth maps in
Fig.~\ref{fig:recognition}e--h use one run per condition with $N=150$. The
replicated analysis (Extended Data Fig.~\ref{ed:recognition-selection}) uses 30
independent runs per condition with $N=150$ and reports means across runs, with
pointwise 95\% percentile intervals from 2,000 bootstrap resamples of the run
means. A condition is eligible for selection if, for both $\Delta g$ and
$\bar d$, the bootstrap interval of the difference between the second- and
first-half means of the sampling window lies within $(-0.01,0.01)$. Two slowly
relaxing SJ conditions required longer runs: $(k_1,k_2)=(4,5)$ with 3,000 sweeps
(sampled from sweep 2,400) and $(3,5)$ with 50,000 sweeps (sampled from sweep
40,000); all 196 conditions then passed. Among eligible conditions with
$\Delta g\ge0.8$ and $\bar d\le0.1$, we select the smallest $K$, breaking ties
by larger $\Delta g$, then smaller $\bar d$, $k_1$ and $k_2$. This gives the
four architectures listed above. The stars in Fig.~\ref{fig:recognition}e--h
apply the same cutoffs and tie-breaking rule to the single-run maps, without
the drift criterion, and identify the same four architectures. Selection uses
reputation statistics only.
The robustness analysis varies the separation cutoff from 0.70 to 0.90 and the
disagreement cutoff from 0.050 to 0.150 in steps of 0.05 and 0.025. The
illustrative image matrices in Fig.~\ref{fig:recognition}a--d are final states of
single runs with $N=90$ after 1,000 sweeps.

\subsection*{Numerical procedures and figure-specific settings}

All panels use $b/c=5$ and $e_1=e_2=0.02$ unless an axis varies them;
Supplementary Note~12 (Supplementary Table~S11) lists every setting, grid and
tolerance.
$\mathcal R$ and $\mathcal R'$ are evaluated in closed form. Homogeneous-DISC
roots (Fig.~\ref{fig:opportunism}; Extended Data
Figs.~\ref{ed:resident-reputation}a--d, \ref{ed:reputation-planes} and
\ref{ed:opportunism-security}) are located by evaluating $F(r)$ on a fine
uniform grid, refining sign changes by bracketing and classifying roots by the
sign of $F'(r)$. At mixed compositions, Eq.~\eqref{eq:m-relaxation} is
integrated from $r_I=1$ by fourth-order Runge--Kutta steps until
$\max_I|\mathcal R(u_I)-r_I|\le10^{-11}$ ($10^{-12}$ for
Fig.~\ref{fig:strategy-dynamics}). For local invasion (Extended Data
Figs.~\ref{ed:benefit-cost}--\ref{ed:readout-evolution}), the resident
relaxation is followed by step-limited Newton refinement, which keeps the
branch reached from $r=1$, to residual $\le10^{-13}$.

All sampled-cooperation results (Fig.~\ref{fig:memory-allocation}; Extended
Data Figs.~\ref{ed:resident-reputation}e, \ref{ed:benefit-cost},
\ref{ed:depth-evolution} and \ref{ed:readout-evolution}e--h) use one engine.
Selected reputations are tabulated exactly on a triangular lattice with 120
subdivisions and interpolated linearly only inside cells in which no reputation
component varies by more than 0.02; elsewhere, including every cell crossed by
a fold, they are computed exactly at the requested composition. Payoffs and the
replicator field are always computed from these reputations, never
interpolated. Trajectories are integrated with the explicit midpoint rule
($\Delta t=0.2$) until $\rho_{\mathrm{int}}$ changes by at most $10^{-7}$ over
each of three consecutive intervals of 100 time units, or for at most
$4\times10^4$ time units. Halving the step, refining the lattice and tightening
the tolerances changed $\bar\rho$ by at most $1.3\times10^{-5}$ and no outcome
class. Initial compositions are 512 points of a randomly shifted Halton
sequence (bases 2 and 3; seed 2026) or, for Extended Data
Fig.~\ref{ed:benefit-cost}, eight independently scrambled 512-point Sobol
sequences~\cite{virtanen2020} (seeds 2026--2033), mapped area-uniformly onto the
simplex by reflecting points of the unit square with $x+y>1$. In Extended Data
Fig.~\ref{ed:benefit-cost}, error bars are standard errors across the eight
samples and therefore describe the sampling of initial compositions; the
dynamics are deterministic. Extended Data Fig.~\ref{ed:resident-reputation}e
initializes reputations once, at $t=0$, as all-good or all-bad and then follows
them quasi-statically along each trajectory by warm-started relaxation, so that
a branch is left only where it ends.

In Fig.~\ref{fig:strategy-dynamics} and Extended Data
Fig.~\ref{ed:readout-evolution}a--d, reputations are relaxed exactly at every
point of an $84\times74$ grid clipped to the simplex. Vertex stability follows
from the two rare-strategy growth rates; edge equilibria are located on
501-point edge scans, and their markers refer to stability along the edge;
interior equilibria are classified by their Jacobian. Under SC without inertia,
all three payoffs are equal on the line
$f_{\mathrm{DISC}}=c/[b(\epsilon-e_2)]\approx0.2126$ (Supplementary Note~6).
Extended Data Fig.~\ref{ed:reputation-planes} reports the highest stable
homogeneous-DISC root and marks bistability (one-step readout plane,
$161\times161$; depth plane, $0\le k_1,k_2\le8$); for every root in the depth
planes, stability in Eq.~\eqref{eq:m-relaxation} agreed with the linearization of
the full latent-state dynamics. On the ALLD--DISC edge (Extended Data
Fig.~\ref{ed:edge-geometry}), Eq.~\eqref{eq:m-self-consistency} reduces exactly
to one equation in $\bar r$, which is solved for 1,071 values of
$f_{\mathrm{DISC}}\in[0.01,1]$. Local invasion maps use $241\times241$ points of
$e_1,e_2\in[0,0.45]$ (Extended Data Fig.~\ref{ed:error-invasion}) and
$801\times801$ points of $(\alpha,\beta)$ (Extended Data
Fig.~\ref{ed:readout-evolution}i--l). In Fig.~\ref{fig:opportunism},
$\Lambda(\delta)$ is evaluated at 5,015 values of $\delta$, concentrated near
$\delta=0$, with sign changes refined by Brent's method; $\delta_{\mathrm c}$ is
the upper end of the interval with $\Lambda>0$. Residents of the alternative
filters are obtained from $r=1$ by damped iteration,
$r\leftarrow0.8r+0.2\,\mathcal R_{\mathrm X}[u_{\mathrm R}(r)]$ with
$\mathcal R_{\mathrm X}\in\{\mathcal R_{\mathrm{agg}},\mathcal R_{\mathrm{score}}\}$.
For Extended Data Fig.~\ref{ed:opportunism-security} ($\ge601$ values of
$\alpha$, 161 values of $b/c\in[2,8]$, $\beta=1$), $\min_\delta H(\delta)$ is
found by comparing $H$ at $\delta=0$, $\delta=1$ and all real roots in
$(10^{-8},1)$ of the numerator of the rational function
$\mathrm dH/\mathrm d\delta$, and a readout is classified as secure if
$1-(b/c)\min_\delta H\le10^{-10}$. Fold endpoints are excluded; endpoints at the
domain boundary or at neutral invasion are included.

\subsection*{Data and code availability}

The code, numerical data and figure scripts used in this study are provided
with the paper (verified computational package; one notebook per figure, a C
engine for the strategy dynamics, input configurations, validation records and
checksums). \textit{[Repository and DOI to be added before submission.]}
 in place of the
% Methods section of Manuscript_formal_20261028.tex, or compiled
% stand-alone through Methods.tex. Bibliography keys used here are
% listed in Methods_references.tex.
% Figure labels follow Manuscript_formal_20261028.tex.
% Notation follows Supplementary Note 1 (single notation system).
% =====================================================================
\section*{Methods}
\label{sec:methods}

\subsection*{Donation game, strategies and private assessment}

We consider an infinitely large, well-mixed population in which randomly
chosen donors and recipients play a donation game. A donor either cooperates,
paying a cost $c>0$ to give the recipient a benefit $b>c$, or defects, which has
no payoff consequences. Individuals use one of three strategies
$I\in\{\mathrm{ALLC},\mathrm{ALLD},\mathrm{DISC}\}$ with frequencies $f_I$,
$\sum_I f_I=1$. ALLC intends to cooperate with every recipient, ALLD intends to
defect against every recipient, and DISC intends to cooperate if and only if it
regards the recipient as good. All individuals share one social norm and one
reputation mechanism (the depths $k_1,k_2$ and readout parameters
$\alpha,\beta$ defined below); strategy evolution changes only the frequencies
$f_I$.

Reputations are private~\cite{okada2018,hilbe2018,radzvilavicius2019}: every
individual keeps its own record of every other individual. In each
interaction, an observer drawn independently of donor and recipient assesses
the donor. A DISC donor chooses its action using its own image of the
recipient, whereas the observer judges the donor's realized action using the
observer's own image of the recipient, so the two can disagree. We use
second-order norms~\cite{ohtsuki2004}. Cooperation with a recipient regarded as
good is assessed positively and defection against such a recipient negatively;
towards a recipient regarded as bad, cooperation and defection are assessed
positively with probabilities $p$ and $q$, respectively. Scoring (SC), Shunning
(SH), Stern Judging (SJ) and Simple Standing (SS) correspond to
$(p,q)=(1,0)$, $(0,0)$, $(0,1)$ and $(1,1)$.

Intended cooperation fails and becomes defection with probability $e_1$;
intended defection is always implemented. After applying the norm, the
observer records the opposite assessment with probability $e_2$. The two errors
are independent, and benefits and costs accrue only for realized cooperation.
Let $P^+_{A\mid Y}$ denote the probability that the recorded assessment is
positive, given the donor's intended action $A\in\{\mathrm C,\mathrm D\}$ and
the observer's image $Y\in\{\mathrm G,\mathrm B\}$ of the recipient. With
$\epsilon=(1-e_1)(1-e_2)+e_1e_2$,
\begin{equation}
\begin{aligned}
P^+_{\mathrm C\mid\mathrm G}&=\epsilon, &
P^+_{\mathrm C\mid\mathrm B}&=e_2+(1-2e_2)\bigl[(1-e_1)p+e_1q\bigr],\\
P^+_{\mathrm D\mid\mathrm G}&=e_2, &
P^+_{\mathrm D\mid\mathrm B}&=e_2+(1-2e_2)q
\end{aligned}
\label{eq:m-assessment}
\end{equation}
(Supplementary Note~2). Implementation error is already averaged into
Eq.~\eqref{eq:m-assessment}. If the observer regards the recipient as good with
probability $r$, intended cooperation and intended defection are therefore
assessed positively with probabilities
\begin{equation}
\Phi_{\mathrm C}(r)=rP^+_{\mathrm C\mid\mathrm G}+(1-r)P^+_{\mathrm C\mid\mathrm B},
\qquad
\Phi_{\mathrm D}(r)=rP^+_{\mathrm D\mid\mathrm G}+(1-r)P^+_{\mathrm D\mid\mathrm B},
\label{eq:m-action-assessment}
\end{equation}
and their difference, the assessment contrast, is
$\Delta\Phi(r)=\Phi_{\mathrm C}(r)-\Phi_{\mathrm D}(r)=\kappa\,[r+(p-q)(1-r)]$
with $\kappa=(1-e_1)(1-2e_2)$. Unless stated otherwise, $b=5$, $c=1$ and
$e_1=e_2=0.02$.

\subsection*{Inertial reputation records and binary readout}

Observer $i$ holds a latent record $L_{ij}$ of target $j$ in the state space
\begin{equation}
\mathcal L=\{\mathrm G,\mathrm W_1,\dots,\mathrm W_{k_1},\mathrm B,
\mathrm S_1,\dots,\mathrm S_{k_2}\},\qquad k_1,k_2\in\{0,1,2,\dots\}.
\label{eq:m-states}
\end{equation}
After each interaction, only the observer's record of the donor changes, and
only through the recorded assessment. A negative assessment moves
$\mathrm G\to\mathrm W_1\to\dots\to\mathrm W_{k_1}\to\mathrm B$ and leaves
$\mathrm B$ unchanged; a positive assessment moves
$\mathrm B\to\mathrm S_1\to\dots\to\mathrm S_{k_2}\to\mathrm G$ and leaves
$\mathrm G$ unchanged. A positive assessment returns every warning state
$\mathrm W_m$ directly to $\mathrm G$, and a negative assessment returns every
recovery state $\mathrm S_n$ directly to $\mathrm B$ (direct reset). If
$k_1=0$, a negative assessment moves $\mathrm G$ directly to $\mathrm B$; if
$k_2=0$, a positive assessment moves $\mathrm B$ directly to $\mathrm G$.
Losing good standing therefore requires $k_1+1$, and regaining it $k_2+1$,
consecutive assessments of the same target by the same observer. We call $k_1$
the tolerance depth, $k_2$ the strictness depth and $K=k_1+k_2$ the total
depth; a record has $K+2$ states (complete transition table in Supplementary Note~3 and
Supplementary Table~S5).

Actions and assessments use a binary image that is sampled from the record
whenever the record is consulted. The image is good with probability
\begin{equation}
\sigma(\mathrm G)=1,\qquad \sigma(\mathrm W_m)=\alpha^m,\qquad
\sigma(\mathrm B)=0,\qquad \sigma(\mathrm S_n)=1-\beta^n,
\qquad 0\le\alpha,\beta\le1 .
\label{eq:m-readout}
\end{equation}
Consultations are independent given the current state and do not change it.
The tolerance strength $\alpha$ sets how often a warning state is still treated
as good, and the strictness strength $\beta$ how often a recovery state is
still treated as bad. We compare three settings: no inertia,
$(k_1,k_2)=(0,0)$; full one-step inertia, $(k_1,k_2)=(1,1)$ with
$\alpha=\beta=1$; and the selected multi-step architectures with full readout,
$(2,0)$, $(3,0)$, $(4,3)$ and $(1,2)$ for SC, SH, SJ and SS (selected from
simulated reputation statistics; see below). Setting $\alpha=\beta=0$ at any
depths also yields instantaneous updating, because the image then always
equals the last recorded assessment. At one-step depths, $(\alpha,\beta)=(1,0)$
and $(0,1)$ give tolerance-only and strictness-only readout.

\subsection*{Stationary reputation map}

Consider one record whose successive recorded assessments are independent and
positive with probability $u$. Its state then follows a Markov chain with
transition matrix $T(u)=uT^{+}+(1-u)T^{-}$, where $T^{+}$ and $T^{-}$ encode the
deterministic updates after positive and negative assessments. For $0<u<1$ the
chain is irreducible and aperiodic. Write $v=1-u$. Because $\mathrm W_m$ can be
entered only from $\mathrm W_{m-1}$ (with $\mathrm W_0\equiv\mathrm G$) after a
negative assessment, and $\mathrm S_n$ only from $\mathrm S_{n-1}$ (with
$\mathrm S_0\equiv\mathrm B$) after a positive one, the stationary distribution
satisfies $\pi_{\mathrm W_m}=v^m\pi_{\mathrm G}$ and
$\pi_{\mathrm S_n}=u^n\pi_{\mathrm B}$. The only transitions from the
good-standing states $\{\mathrm G,\mathrm W_m\}$ to the bad-standing states
$\{\mathrm B,\mathrm S_n\}$ complete a run of $k_1+1$ negative assessments, and
the only reverse transitions complete a run of $k_2+1$ positive assessments.
Balancing these two fluxes gives $v^{k_1+1}\pi_{\mathrm G}=u^{k_2+1}\pi_{\mathrm B}$,
and normalization yields
\begin{equation}
\pi_{\mathrm G}=\frac{u^{k_2+1}}{Z(u)},\quad
\pi_{\mathrm W_m}=\frac{v^{m}u^{k_2+1}}{Z(u)},\quad
\pi_{\mathrm B}=\frac{v^{k_1+1}}{Z(u)},\quad
\pi_{\mathrm S_n}=\frac{u^{n}v^{k_1+1}}{Z(u)},\qquad
Z(u)=v^{k_1}+u^{k_2}-v^{k_1}u^{k_2}
\label{eq:m-stationary}
\end{equation}
(Supplementary Note~4). Averaging the readout probabilities over this
distribution defines the reputation map, the stationary probability that the
target is regarded as good when it is assessed positively with probability $u$:
\begin{equation}
\mathcal R(u)\equiv\mathcal R(u;\alpha,\beta,k_1,k_2)
=\sum_{\ell\in\mathcal L}\pi_\ell(u)\,\sigma(\ell)
=\frac{u^{k_2+1}\sum_{m=0}^{k_1}(\alpha v)^m
+v^{k_1+1}\sum_{n=1}^{k_2}u^n\bigl(1-\beta^n\bigr)}
{v^{k_1}+u^{k_2}-v^{k_1}u^{k_2}} .
\label{eq:m-reputation-map}
\end{equation}
Empty sums are zero, and $\mathcal R(0)=0$ and $\mathcal R(1)=1$ by continuity.
Without inertia, $\mathcal R(u)=u$. For one-step inertia,
$\mathcal R(u)=\{u^2[1+\alpha(1-u)]+(1-\beta)u(1-u)^2\}/(1-u+u^2)$, and with
full readout, $\mathcal R(u)=u^{k_2}(1-v^{k_1+1})/Z(u)$. The map increases with
$\alpha$ and decreases with $\beta$. Two properties are used repeatedly below:
for $k_1\ge1$, $\mathcal R'(1)=1-\alpha$, where the prime denotes $\mathrm d/\mathrm du$;
and for $k_2=0$, $\mathcal R$ is concave on $[0,1]$
(Supplementary Note~4).

\subsection*{Self-consistent private reputations}

Let $r_I$ be the probability that a randomly chosen observer regards a randomly
chosen type-$I$ individual as good, and $\bar r=\sum_I f_I r_I$. We close the
private-assessment problem with two mean-field assumptions, as in previous
analyses of private assessment~\cite{okada2018,radzvilavicius2019,kessinger2023,michelmata2024,kawakatsu2024}.
First, the donor's and the observer's images of a recipient are independent
given the recipient's strategy. Both then regard a random recipient as good
with probability $h_{\mathrm{GG}}=\sum_J f_J r_J^2$, both regard it as bad with
probability $h_{\mathrm{BB}}=\sum_J f_J(1-r_J)^2$, and each of the two ordered
disagreements occurs with probability $h_{\mathrm{GB}}=\sum_J f_J r_J(1-r_J)$.
A type-$I$ donor is therefore assessed positively with probability
\begin{equation}
\begin{aligned}
u_{\mathrm{ALLC}}&=\Phi_{\mathrm C}(\bar r),\qquad
u_{\mathrm{ALLD}}=\Phi_{\mathrm D}(\bar r),\\
u_{\mathrm{DISC}}&=h_{\mathrm{GG}}P^+_{\mathrm C\mid\mathrm G}
+h_{\mathrm{GB}}\bigl(P^+_{\mathrm C\mid\mathrm B}+P^+_{\mathrm D\mid\mathrm G}\bigr)
+h_{\mathrm{BB}}P^+_{\mathrm D\mid\mathrm B}.
\end{aligned}
\label{eq:m-strategy-assessment}
\end{equation}
Second, the assessments entering a record of a type-$I$ individual are treated
as independent with this common probability $u_I$, so that each record has
the stationary response of Eq.~\eqref{eq:m-reputation-map}. Stationary
reputations then satisfy the self-consistency equations
\begin{equation}
r_I^*=\mathcal R\bigl(u_I(\boldsymbol r^*;\boldsymbol f)\bigr),
\qquad I\in\{\mathrm{ALLC},\mathrm{ALLD},\mathrm{DISC}\}.
\label{eq:m-self-consistency}
\end{equation}
Equation~\eqref{eq:m-self-consistency} can have several solutions. Consistent
with records that start in good standing, we select at every composition the
solution reached by the auxiliary relaxation
\begin{equation}
\frac{\mathrm d r_I}{\mathrm d\tau}=\mathcal R(u_I)-r_I,\qquad r_I(0)=1 .
\label{eq:m-relaxation}
\end{equation}
The variable $\tau$ is a computational time, unrelated to evolutionary time
$t$. The fixed points of Eq.~\eqref{eq:m-relaxation} coincide with the
stationary states of the mean-field dynamics of the full distribution of
latent states, $\mathrm d\boldsymbol\pi_I/\mathrm dt'=\boldsymbol\pi_I[T(u_I)-\mathbf I]$
with $r_I=\boldsymbol\pi_I\boldsymbol\sigma$ (Supplementary Note~5). Stability
statements refer to Eq.~\eqref{eq:m-relaxation}; where checked, they agreed
with the linearization of the full latent-state system.

In a homogeneous DISC population with reputation $r$, a resident donor regards
its recipient as good with probability $r$, independently of the observer, and
Eq.~\eqref{eq:m-strategy-assessment} reduces to
\begin{equation}
u_{\mathrm R}(r)=r\,\Phi_{\mathrm C}(r)+(1-r)\,\Phi_{\mathrm D}(r).
\label{eq:m-resident-assessment}
\end{equation}
Resident reputations are the roots of
$F(r)\equiv\mathcal R[u_{\mathrm R}(r)]-r$, and their stability in
Eq.~\eqref{eq:m-relaxation} is given by the sign of $F'(r)$. A fold of a
resident branch satisfies $F(r)=F'(r)=0$, that is
$\mathcal R'[u_{\mathrm R}(r)]\,u_{\mathrm R}'(r)=1$. We denote by $r_{\mathrm R}^*$
the highest strictly stable root, which coincides with the root reached by
relaxation from $r=1$ unless the highest root is a fold, and write
$u_{\mathrm R}^*=u_{\mathrm R}(r_{\mathrm R}^*)$,
$\Phi_A^*=\Phi_A(r_{\mathrm R}^*)$ and
$\Delta\Phi^*=\Delta\Phi(r_{\mathrm R}^*)$. Resident reputations do not depend
on $b/c$.

\subsection*{Strategy evolution and cooperation}

Reputations are assumed to equilibrate quickly relative to strategy change. At
each strategy composition, payoffs are evaluated with the selected stationary
reputations. Per unit time, in which each individual acts on average once as
donor and once as recipient, the expected payoffs are
\begin{equation}
\begin{aligned}
\Pi_{\mathrm{ALLC}}&=(1-e_1)\bigl[b\,(f_{\mathrm{ALLC}}+f_{\mathrm{DISC}}r^*_{\mathrm{ALLC}})-c\bigr],\\
\Pi_{\mathrm{ALLD}}&=(1-e_1)\,b\,(f_{\mathrm{ALLC}}+f_{\mathrm{DISC}}r^*_{\mathrm{ALLD}}),\\
\Pi_{\mathrm{DISC}}&=(1-e_1)\bigl[b\,(f_{\mathrm{ALLC}}+f_{\mathrm{DISC}}r^*_{\mathrm{DISC}})-c\,\bar r^*\bigr],
\end{aligned}
\label{eq:m-payoffs}
\end{equation}
and strategy frequencies follow the replicator dynamics~\cite{hofbauer1998}
\begin{equation}
\frac{\mathrm d f_I}{\mathrm dt}=f_I\bigl(\Pi_I-\bar\Pi\bigr),\qquad
\bar\Pi=\sum_J f_J\Pi_J .
\label{eq:m-replicator}
\end{equation}
The intended and realized cooperation rates are
$\rho_{\mathrm{int}}=f_{\mathrm{ALLC}}+f_{\mathrm{DISC}}\bar r^*$ and
$\rho=(1-e_1)\rho_{\mathrm{int}}$, and $\bar\Pi=(b-c)\rho$. In a homogeneous DISC
population, $\rho_{\mathrm{int}}=r^*_{\mathrm R}$; realized cooperation
cannot exceed $1-e_1=0.98$.

To compare parameter settings globally, we integrate Eq.~\eqref{eq:m-replicator}
from $n_{\mathrm{init}}=512$ initial compositions sampled uniformly by area from
the simplex, using the same compositions for all settings that are compared.
In every scan, each trajectory ended either in the zero-cooperation set (the
ALLD vertex together with the low-reputation segment of the ALLD--DISC edge,
along which the dynamics drift slowly towards ALLD; final $\rho\le7\times10^{-4}$)
or at an attractor with positive cooperation (final $\rho\ge0.34$, and
$0.974$--$0.980$ in nearly all conditions). We therefore define the cooperative
basin as the set of initial compositions whose trajectories reach an attractor
with positive cooperation, and classify a trajectory as cooperative if its
final $\rho\ge0.01$; any threshold inside the observed gap gives the same
classification, and the gap was checked for every condition. We report the
mean long-run cooperation $\bar\rho$ (the average of the final $\rho$ over
initial compositions), the cooperative-basin size $\omega_{\mathrm C}$ (the
fraction of cooperative trajectories) and the cooperation level
$\bar\rho_{\mathrm C}$ within the basin, with
$\bar\rho\approx\omega_{\mathrm C}\bar\rho_{\mathrm C}$. For the total-depth
comparison, the envelope at total depth $K$ is the largest $\bar\rho$ among the
$K+1$ allocations with $k_1+k_2=K$.

\subsection*{Invasion of a discriminator population}

We consider a homogeneous DISC population at $r^*_{\mathrm R}$ and a rare
mutant $\mathrm X$ that uses the residents' norm and record mechanism. Because
the mutant is rare, it interacts with residents only and leaves their
reputations and assessments unchanged. Suppose that $\mathrm X$ intends to help
a resident recipient with probability $\chi_{\mathrm X}$, independently of the
observer's image of that recipient, and acts independently across interactions.
It is then assessed positively with probability
$u_{\mathrm X}=\chi_{\mathrm X}\Phi^*_{\mathrm C}+(1-\chi_{\mathrm X})\Phi^*_{\mathrm D}$,
its stationary reputation is $r^*_{\mathrm X}=\mathcal R(u_{\mathrm X})$, and its
normalized payoff advantage over residents is (Supplementary Note~7)
\begin{equation}
\Lambda_{\mathrm X}\equiv\frac{\Pi_{\mathrm X}-\Pi_{\mathrm R}}{c(1-e_1)}
=\bigl(r^*_{\mathrm R}-\chi_{\mathrm X}\bigr)
-\frac{b}{c}\bigl(r^*_{\mathrm R}-r^*_{\mathrm X}\bigr),
\label{eq:m-mutant-identity}
\end{equation}
the cooperation cost saved minus the benefit lost through reputation. The
resident itself corresponds to $\chi=r^*_{\mathrm R}$ and
$u=u^*_{\mathrm R}=r^*_{\mathrm R}\Phi^*_{\mathrm C}+(1-r^*_{\mathrm R})\Phi^*_{\mathrm D}$.

\textit{Unconditional mutants.} ALLD ($\chi=0$) has reputation
$\mathcal R(\Phi^*_{\mathrm D})$ and ALLC ($\chi=1$) has reputation
$\mathcal R(\Phi^*_{\mathrm C})$, where
$\Phi^*_{\mathrm D}=r^*_{\mathrm R}P^+_{\mathrm D\mid\mathrm G}+(1-r^*_{\mathrm R})P^+_{\mathrm D\mid\mathrm B}$
and $\Phi^*_{\mathrm C}$ analogously. When ALLD is regarded worse and ALLC
better than the residents, Eq.~\eqref{eq:m-mutant-identity} implies that the
discriminator population resists both mutants if and only if
\begin{equation}
R_{\mathrm L}\equiv
\frac{r^*_{\mathrm R}}
{r^*_{\mathrm R}-\mathcal R\bigl[r^*_{\mathrm R}P^+_{\mathrm D\mid\mathrm G}
+(1-r^*_{\mathrm R})P^+_{\mathrm D\mid\mathrm B}\bigr]}
<\frac bc<
\frac{1-r^*_{\mathrm R}}
{\mathcal R\bigl[r^*_{\mathrm R}P^+_{\mathrm C\mid\mathrm G}
+(1-r^*_{\mathrm R})P^+_{\mathrm C\mid\mathrm B}\bigr]-r^*_{\mathrm R}}
\equiv R_{\mathrm U}.
\label{eq:m-window}
\end{equation}
If ALLD is not regarded worse than the residents, no $b/c$ excludes it; if ALLC
is not regarded better, it cannot invade. Because
$u^*_{\mathrm R}=r^*_{\mathrm R}\Phi^*_{\mathrm C}+(1-r^*_{\mathrm R})\Phi^*_{\mathrm D}$,
the interval in Eq.~\eqref{eq:m-window} is non-empty if and only if
$r^*_{\mathrm R}\mathcal R(\Phi^*_{\mathrm C})+(1-r^*_{\mathrm R})\mathcal R(\Phi^*_{\mathrm D})
<\mathcal R(u^*_{\mathrm R})$, that is, if $\mathcal R$ at the resident's
assessment probability lies above the chord joining its values at
$\Phi^*_{\mathrm D}$ and $\Phi^*_{\mathrm C}$. Without inertia,
$\mathcal R$ is linear and $R_{\mathrm L}=R_{\mathrm U}=1/\Delta\Phi^*$, so no
benefit-to-cost ratio protects a discriminator population (Supplementary
Note~7). Numerically, we classify local outcomes from the signs of the
unnormalized payoff differences
$\lambda_{\mathrm{ALLD}}=\Pi_{\mathrm{ALLD}}-\Pi_{\mathrm R}$ and
$\lambda_{\mathrm{ALLC}}=\Pi_{\mathrm{ALLC}}-\Pi_{\mathrm R}$ (with $c=1$), which also
covers cases without finite thresholds: ALLD is not excluded if
$\lambda_{\mathrm{ALLD}}\ge-10^{-10}$; otherwise ALLC invades if
$\lambda_{\mathrm{ALLC}}>10^{-10}$; otherwise the residents resist both.

\textit{Opportunistic discriminators.} An opportunistic discriminator (O-DISC)
behaves like DISC, except that each intended cooperation towards a recipient it
regards as good is replaced, independently and before implementation error, by
intended defection with probability $\delta\in[0,1]$. Thus
$\chi=(1-\delta)r^*_{\mathrm R}$, and $\delta=0$ and $\delta=1$ recover DISC and
ALLD. The mutant's positive-assessment probability and reputation are
\begin{equation}
u_{\mathrm X}(\delta)=u^*_{\mathrm R}-\delta\,r^*_{\mathrm R}\,\Delta\Phi^*,
\qquad
r^*_{\mathrm X}(\delta)=\mathcal R\bigl[u_{\mathrm X}(\delta)\bigr],
\label{eq:m-odisc-reputation}
\end{equation}
and Eq.~\eqref{eq:m-mutant-identity} gives the normalized invasion payoff and
its slope at the resident,
\begin{equation}
\Lambda(\delta)=\delta\,r^*_{\mathrm R}
-\frac{b}{c}\bigl[r^*_{\mathrm R}-r^*_{\mathrm X}(\delta)\bigr],
\qquad
\Lambda'(0)=r^*_{\mathrm R}\Bigl[1-\frac{b}{c}\,\Delta\Phi^*\,
\mathcal R'(u^*_{\mathrm R})\Bigr].
\label{eq:m-odisc-invasion}
\end{equation}
A positive $\Lambda(\delta)$ means that the mutant can invade; $\Lambda(0)=0$,
and $\Lambda(1)$ is the normalized ALLD advantage. For $r^*_{\mathrm R}>0$,
$\Lambda(\delta)=\delta r^*_{\mathrm R}[1-(b/c)H(\delta)]$, where
$H(\delta)=[r^*_{\mathrm R}-r^*_{\mathrm X}(\delta)]/(\delta r^*_{\mathrm R})$ is the
reputation lost per unit of withheld cooperation, extended continuously by
$H(0)=\Delta\Phi^*\mathcal R'(u^*_{\mathrm R})$. The residents therefore exclude
every O-DISC mutant if and only if
\begin{equation}
\min_{0\le\delta\le1}\frac{r^*_{\mathrm R}-r^*_{\mathrm X}(\delta)}{\delta\,r^*_{\mathrm R}}
\;\ge\;\frac{c}{b},
\label{eq:m-security}
\end{equation}
with equality permitting neutral mutants. Because $H(\delta)$ equals
$\Delta\Phi^*$ times the slope of the chord of $\mathcal R$ between
$u_{\mathrm X}(\delta)$ and $u^*_{\mathrm R}$, for $\Delta\Phi^*>0$ the minimum
is attained at $\delta\to0$ whenever $\mathcal R$ is concave on that interval
(always for $k_2=0$), and Eq.~\eqref{eq:m-security} then reduces to
$\Lambda'(0)\le0$; in general, the whole interval must be checked. Near an
all-good resident with $k_1\ge1$, $\mathcal R'(u^*_{\mathrm R})\approx1-\alpha$,
so full tolerance readout makes weak opportunism cheap. In the error-free limit with $\beta=1$ and $k_1\ge1$, an
all-good resident excludes every O-DISC mutant if and only if
$\alpha\le1-c/b$, for every norm and every $k_2$ (Supplementary Note~8); at
$b/c=5$ this bound is the readout $\alpha=0.8$ used in Fig.~\ref{fig:opportunism}
and Extended Data Fig.~\ref{ed:opportunism-security}.

For fixed norm, depths, $\beta$ and errors, a readout $\alpha$ is secure if a
strictly stable high-reputation resident exists and satisfies
Eq.~\eqref{eq:m-security}. The lower boundary $\alpha_{\mathrm L}$ of a secure
interval, when it is not $\alpha=0$, is a fold of the high-reputation branch,
$F(r,\alpha)=\partial_rF(r,\alpha)=0$ with
$F(r,\alpha)=\mathcal R[u_{\mathrm R}(r);\alpha,\beta,k_1,k_2]-r$; the fold
itself is excluded and does not depend on $b/c$. The upper boundary
$\alpha_{\mathrm U}$, when it is not $\alpha=1$, is where
Eq.~\eqref{eq:m-security} holds with equality. In every case we computed it was
determined by $\delta\to0$, that is by
$(b/c)\,\Delta\Phi(r)\,\mathcal R'[u_{\mathrm R}(r);\alpha_{\mathrm U}]=1$
together with $F(r,\alpha_{\mathrm U})=0$. When the fold and the local condition
meet, the secure interval closes; this happens at
$b/c=u'_{\mathrm R}(r)/\Delta\Phi(r)$ evaluated at the fold, which equals 2 for
Shunning (Supplementary Note~8).

\textit{Alternative reputation filters.} For Fig.~\ref{fig:opportunism}a we
replace $\mathcal R$ by the stationary response of two alternative mechanisms
while keeping the norm, errors, resident selection and
Eqs.~\eqref{eq:m-odisc-reputation}--\eqref{eq:m-odisc-invasion} unchanged.
Aggregation, adapted from ref.~\cite{michelmata2024}, assigns a good image if at
least $Q$ of $M$ independently sampled assessments are positive,
$\mathcal R_{\mathrm{agg}}(u)=\sum_{j=Q}^{M}\binom Mj u^j(1-u)^{M-j}$, with
$(M,Q)=(2,1)$, so that $\mathcal R_{\mathrm{agg}}(u)=2u-u^2$. Quantitative
assessment, adapted from ref.~\cite{schmid2023}, keeps an integer score
$z\in\{-z_{\max},\dots,z_{\max}\}$ that moves up or down by one after a positive
or negative assessment (attempts to leave the range leave the score unchanged)
and is read as good if $z\ge z_0$; with $t=u/(1-u)$,
$\mathcal R_{\mathrm{score}}(u)=\sum_{j=z_{\max}+z_0}^{2z_{\max}}t^j\big/\sum_{j=0}^{2z_{\max}}t^j$,
and we use $z_{\max}=5$, $z_0=0$. Neither mechanism was optimized or matched in
memory to the inertial record.

\subsection*{Finite-population reputation simulations}

Individual-based simulations test reputation formation without the mean-field
closure, at fixed strategy composition; they do not simulate strategy
evolution. A population of $N$ individuals contains fixed numbers $n_I$ of each
strategy. Every DISC individual stores a latent record of every individual,
including itself, and all records start in $\mathrm G$. Records held by ALLC and
ALLD individuals are not stored because they never affect actions or
assessments. In each event: (i) a donor is drawn uniformly from the population
and a recipient uniformly from the remaining $N-1$ individuals, and an observer
is drawn independently and uniformly from the DISC individuals (the observer
may coincide with donor or recipient); (ii) the donor forms its intended
action, a DISC donor by sampling its own image of the recipient from
Eq.~\eqref{eq:m-readout}; (iii) implementation error is applied; (iv) the
observer independently samples its own image of the recipient, applies the norm
to the realized action and reverses the assessment with probability $e_2$; and
(v) the observer's record of the donor is updated. A sweep consists of
$n_{\mathrm{DISC}}N$ events, one update per stored record on average.

At sampled times we use $s_{ij}=\sigma(L_{ij})$, which averages exactly over
readout randomness. For target $j$, the mean goodness and the probability that
two distinct DISC observers give different images are
\begin{equation}
g_j=\frac{1}{n_{\mathrm{DISC}}}\sum_i s_{ij},\qquad
d_j=\binom{n_{\mathrm{DISC}}}{2}^{-1}\sum_{i<i'}
\bigl[s_{ij}(1-s_{i'j})+(1-s_{ij})s_{i'j}\bigr],
\label{eq:m-finite-estimators}
\end{equation}
where $i,i'$ run over DISC observers. Averaging over type-$I$ targets gives
$g_I$ and $d_I$, which are then averaged over sampling times. The reputational
separation and mean disagreement are
\begin{equation}
\Delta g=\frac{g_{\mathrm{DISC}}+g_{\mathrm{ALLC}}}{2}-g_{\mathrm{ALLD}},
\qquad
\bar d=\sum_I\frac{n_I}{N}\,d_I .
\label{eq:m-recognition}
\end{equation}
Unlike the mean-field theory, the simulations retain target-specific variation
and correlations between the private records of different observers.

The recognition analyses use equal strategy numbers, $\alpha=\beta=1$ and
$0\le k_1,k_2\le6$. Each run lasts 1,000 sweeps and is sampled every five sweeps
from sweep 800 to sweep 1,000 (41 samples). The depth maps in
Fig.~\ref{fig:recognition}e--h use one run per condition with $N=150$. The
replicated analysis (Extended Data Fig.~\ref{ed:recognition-selection}) uses 30
independent runs per condition with $N=150$ and reports means across runs, with
pointwise 95\% percentile intervals from 2,000 bootstrap resamples of the run
means. A condition is eligible for selection if, for both $\Delta g$ and
$\bar d$, the bootstrap interval of the difference between the second- and
first-half means of the sampling window lies within $(-0.01,0.01)$. Two slowly
relaxing SJ conditions required longer runs: $(k_1,k_2)=(4,5)$ with 3,000 sweeps
(sampled from sweep 2,400) and $(3,5)$ with 50,000 sweeps (sampled from sweep
40,000); all 196 conditions then passed. Among eligible conditions with
$\Delta g\ge0.8$ and $\bar d\le0.1$, we select the smallest $K$, breaking ties
by larger $\Delta g$, then smaller $\bar d$, $k_1$ and $k_2$. This gives the
four architectures listed above. The stars in Fig.~\ref{fig:recognition}e--h
apply the same cutoffs and tie-breaking rule to the single-run maps, without
the drift criterion, and identify the same four architectures. Selection uses
reputation statistics only.
The robustness analysis varies the separation cutoff from 0.70 to 0.90 and the
disagreement cutoff from 0.050 to 0.150 in steps of 0.05 and 0.025. The
illustrative image matrices in Fig.~\ref{fig:recognition}a--d are final states of
single runs with $N=90$ after 1,000 sweeps.

\subsection*{Numerical procedures and figure-specific settings}

All panels use $b/c=5$ and $e_1=e_2=0.02$ unless an axis varies them;
Supplementary Note~12 (Supplementary Table~S11) lists every setting, grid and
tolerance.
$\mathcal R$ and $\mathcal R'$ are evaluated in closed form. Homogeneous-DISC
roots (Fig.~\ref{fig:opportunism}; Extended Data
Figs.~\ref{ed:resident-reputation}a--d, \ref{ed:reputation-planes} and
\ref{ed:opportunism-security}) are located by evaluating $F(r)$ on a fine
uniform grid, refining sign changes by bracketing and classifying roots by the
sign of $F'(r)$. At mixed compositions, Eq.~\eqref{eq:m-relaxation} is
integrated from $r_I=1$ by fourth-order Runge--Kutta steps until
$\max_I|\mathcal R(u_I)-r_I|\le10^{-11}$ ($10^{-12}$ for
Fig.~\ref{fig:strategy-dynamics}). For local invasion (Extended Data
Figs.~\ref{ed:benefit-cost}--\ref{ed:readout-evolution}), the resident
relaxation is followed by step-limited Newton refinement, which keeps the
branch reached from $r=1$, to residual $\le10^{-13}$.

All sampled-cooperation results (Fig.~\ref{fig:memory-allocation}; Extended
Data Figs.~\ref{ed:resident-reputation}e, \ref{ed:benefit-cost},
\ref{ed:depth-evolution} and \ref{ed:readout-evolution}e--h) use one engine.
Selected reputations are tabulated exactly on a triangular lattice with 120
subdivisions and interpolated linearly only inside cells in which no reputation
component varies by more than 0.02; elsewhere, including every cell crossed by
a fold, they are computed exactly at the requested composition. Payoffs and the
replicator field are always computed from these reputations, never
interpolated. Trajectories are integrated with the explicit midpoint rule
($\Delta t=0.2$) until $\rho_{\mathrm{int}}$ changes by at most $10^{-7}$ over
each of three consecutive intervals of 100 time units, or for at most
$4\times10^4$ time units. Halving the step, refining the lattice and tightening
the tolerances changed $\bar\rho$ by at most $1.3\times10^{-5}$ and no outcome
class. Initial compositions are 512 points of a randomly shifted Halton
sequence (bases 2 and 3; seed 2026) or, for Extended Data
Fig.~\ref{ed:benefit-cost}, eight independently scrambled 512-point Sobol
sequences~\cite{virtanen2020} (seeds 2026--2033), mapped area-uniformly onto the
simplex by reflecting points of the unit square with $x+y>1$. In Extended Data
Fig.~\ref{ed:benefit-cost}, error bars are standard errors across the eight
samples and therefore describe the sampling of initial compositions; the
dynamics are deterministic. Extended Data Fig.~\ref{ed:resident-reputation}e
initializes reputations once, at $t=0$, as all-good or all-bad and then follows
them quasi-statically along each trajectory by warm-started relaxation, so that
a branch is left only where it ends.

In Fig.~\ref{fig:strategy-dynamics} and Extended Data
Fig.~\ref{ed:readout-evolution}a--d, reputations are relaxed exactly at every
point of an $84\times74$ grid clipped to the simplex. Vertex stability follows
from the two rare-strategy growth rates; edge equilibria are located on
501-point edge scans, and their markers refer to stability along the edge;
interior equilibria are classified by their Jacobian. Under SC without inertia,
all three payoffs are equal on the line
$f_{\mathrm{DISC}}=c/[b(\epsilon-e_2)]\approx0.2126$ (Supplementary Note~6).
Extended Data Fig.~\ref{ed:reputation-planes} reports the highest stable
homogeneous-DISC root and marks bistability (one-step readout plane,
$161\times161$; depth plane, $0\le k_1,k_2\le8$); for every root in the depth
planes, stability in Eq.~\eqref{eq:m-relaxation} agreed with the linearization of
the full latent-state dynamics. On the ALLD--DISC edge (Extended Data
Fig.~\ref{ed:edge-geometry}), Eq.~\eqref{eq:m-self-consistency} reduces exactly
to one equation in $\bar r$, which is solved for 1,071 values of
$f_{\mathrm{DISC}}\in[0.01,1]$. Local invasion maps use $241\times241$ points of
$e_1,e_2\in[0,0.45]$ (Extended Data Fig.~\ref{ed:error-invasion}) and
$801\times801$ points of $(\alpha,\beta)$ (Extended Data
Fig.~\ref{ed:readout-evolution}i--l). In Fig.~\ref{fig:opportunism},
$\Lambda(\delta)$ is evaluated at 5,015 values of $\delta$, concentrated near
$\delta=0$, with sign changes refined by Brent's method; $\delta_{\mathrm c}$ is
the upper end of the interval with $\Lambda>0$. Residents of the alternative
filters are obtained from $r=1$ by damped iteration,
$r\leftarrow0.8r+0.2\,\mathcal R_{\mathrm X}[u_{\mathrm R}(r)]$ with
$\mathcal R_{\mathrm X}\in\{\mathcal R_{\mathrm{agg}},\mathcal R_{\mathrm{score}}\}$.
For Extended Data Fig.~\ref{ed:opportunism-security} ($\ge601$ values of
$\alpha$, 161 values of $b/c\in[2,8]$, $\beta=1$), $\min_\delta H(\delta)$ is
found by comparing $H$ at $\delta=0$, $\delta=1$ and all real roots in
$(10^{-8},1)$ of the numerator of the rational function
$\mathrm dH/\mathrm d\delta$, and a readout is classified as secure if
$1-(b/c)\min_\delta H\le10^{-10}$. Fold endpoints are excluded; endpoints at the
domain boundary or at neutral invasion are included.

\subsection*{Data and code availability}

The code, numerical data and figure scripts used in this study are provided
with the paper (verified computational package; one notebook per figure, a C
engine for the strategy dynamics, input configurations, validation records and
checksums). \textit{[Repository and DOI to be added before submission.]}
 in place of the
% Methods section of Manuscript_formal_20261028.tex, or compiled
% stand-alone through Methods.tex. Bibliography keys used here are
% listed in Methods_references.tex.
% Figure labels follow Manuscript_formal_20261028.tex.
% Notation follows Supplementary Note 1 (single notation system).
% =====================================================================
\section*{Methods}
\label{sec:methods}

\subsection*{Donation game, strategies and private assessment}

We consider an infinitely large, well-mixed population in which randomly
chosen donors and recipients play a donation game. A donor either cooperates,
paying a cost $c>0$ to give the recipient a benefit $b>c$, or defects, which has
no payoff consequences. Individuals use one of three strategies
$I\in\{\mathrm{ALLC},\mathrm{ALLD},\mathrm{DISC}\}$ with frequencies $f_I$,
$\sum_I f_I=1$. ALLC intends to cooperate with every recipient, ALLD intends to
defect against every recipient, and DISC intends to cooperate if and only if it
regards the recipient as good. All individuals share one social norm and one
reputation mechanism (the depths $k_1,k_2$ and readout parameters
$\alpha,\beta$ defined below); strategy evolution changes only the frequencies
$f_I$.

Reputations are private~\cite{okada2018,hilbe2018,radzvilavicius2019}: every
individual keeps its own record of every other individual. In each
interaction, an observer drawn independently of donor and recipient assesses
the donor. A DISC donor chooses its action using its own image of the
recipient, whereas the observer judges the donor's realized action using the
observer's own image of the recipient, so the two can disagree. We use
second-order norms~\cite{ohtsuki2004}. Cooperation with a recipient regarded as
good is assessed positively and defection against such a recipient negatively;
towards a recipient regarded as bad, cooperation and defection are assessed
positively with probabilities $p$ and $q$, respectively. Scoring (SC), Shunning
(SH), Stern Judging (SJ) and Simple Standing (SS) correspond to
$(p,q)=(1,0)$, $(0,0)$, $(0,1)$ and $(1,1)$.

Intended cooperation fails and becomes defection with probability $e_1$;
intended defection is always implemented. After applying the norm, the
observer records the opposite assessment with probability $e_2$. The two errors
are independent, and benefits and costs accrue only for realized cooperation.
Let $P^+_{A\mid Y}$ denote the probability that the recorded assessment is
positive, given the donor's intended action $A\in\{\mathrm C,\mathrm D\}$ and
the observer's image $Y\in\{\mathrm G,\mathrm B\}$ of the recipient. With
$\epsilon=(1-e_1)(1-e_2)+e_1e_2$,
\begin{equation}
\begin{aligned}
P^+_{\mathrm C\mid\mathrm G}&=\epsilon, &
P^+_{\mathrm C\mid\mathrm B}&=e_2+(1-2e_2)\bigl[(1-e_1)p+e_1q\bigr],\\
P^+_{\mathrm D\mid\mathrm G}&=e_2, &
P^+_{\mathrm D\mid\mathrm B}&=e_2+(1-2e_2)q
\end{aligned}
\label{eq:m-assessment}
\end{equation}
(Supplementary Note~2). Implementation error is already averaged into
Eq.~\eqref{eq:m-assessment}. If the observer regards the recipient as good with
probability $r$, intended cooperation and intended defection are therefore
assessed positively with probabilities
\begin{equation}
\Phi_{\mathrm C}(r)=rP^+_{\mathrm C\mid\mathrm G}+(1-r)P^+_{\mathrm C\mid\mathrm B},
\qquad
\Phi_{\mathrm D}(r)=rP^+_{\mathrm D\mid\mathrm G}+(1-r)P^+_{\mathrm D\mid\mathrm B},
\label{eq:m-action-assessment}
\end{equation}
and their difference, the assessment contrast, is
$\Delta\Phi(r)=\Phi_{\mathrm C}(r)-\Phi_{\mathrm D}(r)=\kappa\,[r+(p-q)(1-r)]$
with $\kappa=(1-e_1)(1-2e_2)$. Unless stated otherwise, $b=5$, $c=1$ and
$e_1=e_2=0.02$.

\subsection*{Inertial reputation records and binary readout}

Observer $i$ holds a latent record $L_{ij}$ of target $j$ in the state space
\begin{equation}
\mathcal L=\{\mathrm G,\mathrm W_1,\dots,\mathrm W_{k_1},\mathrm B,
\mathrm S_1,\dots,\mathrm S_{k_2}\},\qquad k_1,k_2\in\{0,1,2,\dots\}.
\label{eq:m-states}
\end{equation}
After each interaction, only the observer's record of the donor changes, and
only through the recorded assessment. A negative assessment moves
$\mathrm G\to\mathrm W_1\to\dots\to\mathrm W_{k_1}\to\mathrm B$ and leaves
$\mathrm B$ unchanged; a positive assessment moves
$\mathrm B\to\mathrm S_1\to\dots\to\mathrm S_{k_2}\to\mathrm G$ and leaves
$\mathrm G$ unchanged. A positive assessment returns every warning state
$\mathrm W_m$ directly to $\mathrm G$, and a negative assessment returns every
recovery state $\mathrm S_n$ directly to $\mathrm B$ (direct reset). If
$k_1=0$, a negative assessment moves $\mathrm G$ directly to $\mathrm B$; if
$k_2=0$, a positive assessment moves $\mathrm B$ directly to $\mathrm G$.
Losing good standing therefore requires $k_1+1$, and regaining it $k_2+1$,
consecutive assessments of the same target by the same observer. We call $k_1$
the tolerance depth, $k_2$ the strictness depth and $K=k_1+k_2$ the total
depth; a record has $K+2$ states (complete transition table in Supplementary Note~3 and
Supplementary Table~S5).

Actions and assessments use a binary image that is sampled from the record
whenever the record is consulted. The image is good with probability
\begin{equation}
\sigma(\mathrm G)=1,\qquad \sigma(\mathrm W_m)=\alpha^m,\qquad
\sigma(\mathrm B)=0,\qquad \sigma(\mathrm S_n)=1-\beta^n,
\qquad 0\le\alpha,\beta\le1 .
\label{eq:m-readout}
\end{equation}
Consultations are independent given the current state and do not change it.
The tolerance strength $\alpha$ sets how often a warning state is still treated
as good, and the strictness strength $\beta$ how often a recovery state is
still treated as bad. We compare three settings: no inertia,
$(k_1,k_2)=(0,0)$; full one-step inertia, $(k_1,k_2)=(1,1)$ with
$\alpha=\beta=1$; and the selected multi-step architectures with full readout,
$(2,0)$, $(3,0)$, $(4,3)$ and $(1,2)$ for SC, SH, SJ and SS (selected from
simulated reputation statistics; see below). Setting $\alpha=\beta=0$ at any
depths also yields instantaneous updating, because the image then always
equals the last recorded assessment. At one-step depths, $(\alpha,\beta)=(1,0)$
and $(0,1)$ give tolerance-only and strictness-only readout.

\subsection*{Stationary reputation map}

Consider one record whose successive recorded assessments are independent and
positive with probability $u$. Its state then follows a Markov chain with
transition matrix $T(u)=uT^{+}+(1-u)T^{-}$, where $T^{+}$ and $T^{-}$ encode the
deterministic updates after positive and negative assessments. For $0<u<1$ the
chain is irreducible and aperiodic. Write $v=1-u$. Because $\mathrm W_m$ can be
entered only from $\mathrm W_{m-1}$ (with $\mathrm W_0\equiv\mathrm G$) after a
negative assessment, and $\mathrm S_n$ only from $\mathrm S_{n-1}$ (with
$\mathrm S_0\equiv\mathrm B$) after a positive one, the stationary distribution
satisfies $\pi_{\mathrm W_m}=v^m\pi_{\mathrm G}$ and
$\pi_{\mathrm S_n}=u^n\pi_{\mathrm B}$. The only transitions from the
good-standing states $\{\mathrm G,\mathrm W_m\}$ to the bad-standing states
$\{\mathrm B,\mathrm S_n\}$ complete a run of $k_1+1$ negative assessments, and
the only reverse transitions complete a run of $k_2+1$ positive assessments.
Balancing these two fluxes gives $v^{k_1+1}\pi_{\mathrm G}=u^{k_2+1}\pi_{\mathrm B}$,
and normalization yields
\begin{equation}
\pi_{\mathrm G}=\frac{u^{k_2+1}}{Z(u)},\quad
\pi_{\mathrm W_m}=\frac{v^{m}u^{k_2+1}}{Z(u)},\quad
\pi_{\mathrm B}=\frac{v^{k_1+1}}{Z(u)},\quad
\pi_{\mathrm S_n}=\frac{u^{n}v^{k_1+1}}{Z(u)},\qquad
Z(u)=v^{k_1}+u^{k_2}-v^{k_1}u^{k_2}
\label{eq:m-stationary}
\end{equation}
(Supplementary Note~4). Averaging the readout probabilities over this
distribution defines the reputation map, the stationary probability that the
target is regarded as good when it is assessed positively with probability $u$:
\begin{equation}
\mathcal R(u)\equiv\mathcal R(u;\alpha,\beta,k_1,k_2)
=\sum_{\ell\in\mathcal L}\pi_\ell(u)\,\sigma(\ell)
=\frac{u^{k_2+1}\sum_{m=0}^{k_1}(\alpha v)^m
+v^{k_1+1}\sum_{n=1}^{k_2}u^n\bigl(1-\beta^n\bigr)}
{v^{k_1}+u^{k_2}-v^{k_1}u^{k_2}} .
\label{eq:m-reputation-map}
\end{equation}
Empty sums are zero, and $\mathcal R(0)=0$ and $\mathcal R(1)=1$ by continuity.
Without inertia, $\mathcal R(u)=u$. For one-step inertia,
$\mathcal R(u)=\{u^2[1+\alpha(1-u)]+(1-\beta)u(1-u)^2\}/(1-u+u^2)$, and with
full readout, $\mathcal R(u)=u^{k_2}(1-v^{k_1+1})/Z(u)$. The map increases with
$\alpha$ and decreases with $\beta$. Two properties are used repeatedly below:
for $k_1\ge1$, $\mathcal R'(1)=1-\alpha$, where the prime denotes $\mathrm d/\mathrm du$;
and for $k_2=0$, $\mathcal R$ is concave on $[0,1]$
(Supplementary Note~4).

\subsection*{Self-consistent private reputations}

Let $r_I$ be the probability that a randomly chosen observer regards a randomly
chosen type-$I$ individual as good, and $\bar r=\sum_I f_I r_I$. We close the
private-assessment problem with two mean-field assumptions, as in previous
analyses of private assessment~\cite{okada2018,radzvilavicius2019,kessinger2023,michelmata2024,kawakatsu2024}.
First, the donor's and the observer's images of a recipient are independent
given the recipient's strategy. Both then regard a random recipient as good
with probability $h_{\mathrm{GG}}=\sum_J f_J r_J^2$, both regard it as bad with
probability $h_{\mathrm{BB}}=\sum_J f_J(1-r_J)^2$, and each of the two ordered
disagreements occurs with probability $h_{\mathrm{GB}}=\sum_J f_J r_J(1-r_J)$.
A type-$I$ donor is therefore assessed positively with probability
\begin{equation}
\begin{aligned}
u_{\mathrm{ALLC}}&=\Phi_{\mathrm C}(\bar r),\qquad
u_{\mathrm{ALLD}}=\Phi_{\mathrm D}(\bar r),\\
u_{\mathrm{DISC}}&=h_{\mathrm{GG}}P^+_{\mathrm C\mid\mathrm G}
+h_{\mathrm{GB}}\bigl(P^+_{\mathrm C\mid\mathrm B}+P^+_{\mathrm D\mid\mathrm G}\bigr)
+h_{\mathrm{BB}}P^+_{\mathrm D\mid\mathrm B}.
\end{aligned}
\label{eq:m-strategy-assessment}
\end{equation}
Second, the assessments entering a record of a type-$I$ individual are treated
as independent with this common probability $u_I$, so that each record has
the stationary response of Eq.~\eqref{eq:m-reputation-map}. Stationary
reputations then satisfy the self-consistency equations
\begin{equation}
r_I^*=\mathcal R\bigl(u_I(\boldsymbol r^*;\boldsymbol f)\bigr),
\qquad I\in\{\mathrm{ALLC},\mathrm{ALLD},\mathrm{DISC}\}.
\label{eq:m-self-consistency}
\end{equation}
Equation~\eqref{eq:m-self-consistency} can have several solutions. Consistent
with records that start in good standing, we select at every composition the
solution reached by the auxiliary relaxation
\begin{equation}
\frac{\mathrm d r_I}{\mathrm d\tau}=\mathcal R(u_I)-r_I,\qquad r_I(0)=1 .
\label{eq:m-relaxation}
\end{equation}
The variable $\tau$ is a computational time, unrelated to evolutionary time
$t$. The fixed points of Eq.~\eqref{eq:m-relaxation} coincide with the
stationary states of the mean-field dynamics of the full distribution of
latent states, $\mathrm d\boldsymbol\pi_I/\mathrm dt'=\boldsymbol\pi_I[T(u_I)-\mathbf I]$
with $r_I=\boldsymbol\pi_I\boldsymbol\sigma$ (Supplementary Note~5). Stability
statements refer to Eq.~\eqref{eq:m-relaxation}; where checked, they agreed
with the linearization of the full latent-state system.

In a homogeneous DISC population with reputation $r$, a resident donor regards
its recipient as good with probability $r$, independently of the observer, and
Eq.~\eqref{eq:m-strategy-assessment} reduces to
\begin{equation}
u_{\mathrm R}(r)=r\,\Phi_{\mathrm C}(r)+(1-r)\,\Phi_{\mathrm D}(r).
\label{eq:m-resident-assessment}
\end{equation}
Resident reputations are the roots of
$F(r)\equiv\mathcal R[u_{\mathrm R}(r)]-r$, and their stability in
Eq.~\eqref{eq:m-relaxation} is given by the sign of $F'(r)$. A fold of a
resident branch satisfies $F(r)=F'(r)=0$, that is
$\mathcal R'[u_{\mathrm R}(r)]\,u_{\mathrm R}'(r)=1$. We denote by $r_{\mathrm R}^*$
the highest strictly stable root, which coincides with the root reached by
relaxation from $r=1$ unless the highest root is a fold, and write
$u_{\mathrm R}^*=u_{\mathrm R}(r_{\mathrm R}^*)$,
$\Phi_A^*=\Phi_A(r_{\mathrm R}^*)$ and
$\Delta\Phi^*=\Delta\Phi(r_{\mathrm R}^*)$. Resident reputations do not depend
on $b/c$.

\subsection*{Strategy evolution and cooperation}

Reputations are assumed to equilibrate quickly relative to strategy change. At
each strategy composition, payoffs are evaluated with the selected stationary
reputations. Per unit time, in which each individual acts on average once as
donor and once as recipient, the expected payoffs are
\begin{equation}
\begin{aligned}
\Pi_{\mathrm{ALLC}}&=(1-e_1)\bigl[b\,(f_{\mathrm{ALLC}}+f_{\mathrm{DISC}}r^*_{\mathrm{ALLC}})-c\bigr],\\
\Pi_{\mathrm{ALLD}}&=(1-e_1)\,b\,(f_{\mathrm{ALLC}}+f_{\mathrm{DISC}}r^*_{\mathrm{ALLD}}),\\
\Pi_{\mathrm{DISC}}&=(1-e_1)\bigl[b\,(f_{\mathrm{ALLC}}+f_{\mathrm{DISC}}r^*_{\mathrm{DISC}})-c\,\bar r^*\bigr],
\end{aligned}
\label{eq:m-payoffs}
\end{equation}
and strategy frequencies follow the replicator dynamics~\cite{hofbauer1998}
\begin{equation}
\frac{\mathrm d f_I}{\mathrm dt}=f_I\bigl(\Pi_I-\bar\Pi\bigr),\qquad
\bar\Pi=\sum_J f_J\Pi_J .
\label{eq:m-replicator}
\end{equation}
The intended and realized cooperation rates are
$\rho_{\mathrm{int}}=f_{\mathrm{ALLC}}+f_{\mathrm{DISC}}\bar r^*$ and
$\rho=(1-e_1)\rho_{\mathrm{int}}$, and $\bar\Pi=(b-c)\rho$. In a homogeneous DISC
population, $\rho_{\mathrm{int}}=r^*_{\mathrm R}$; realized cooperation
cannot exceed $1-e_1=0.98$.

To compare parameter settings globally, we integrate Eq.~\eqref{eq:m-replicator}
from $n_{\mathrm{init}}=512$ initial compositions sampled uniformly by area from
the simplex, using the same compositions for all settings that are compared.
In every scan, each trajectory ended either in the zero-cooperation set (the
ALLD vertex together with the low-reputation segment of the ALLD--DISC edge,
along which the dynamics drift slowly towards ALLD; final $\rho\le7\times10^{-4}$)
or at an attractor with positive cooperation (final $\rho\ge0.34$, and
$0.974$--$0.980$ in nearly all conditions). We therefore define the cooperative
basin as the set of initial compositions whose trajectories reach an attractor
with positive cooperation, and classify a trajectory as cooperative if its
final $\rho\ge0.01$; any threshold inside the observed gap gives the same
classification, and the gap was checked for every condition. We report the
mean long-run cooperation $\bar\rho$ (the average of the final $\rho$ over
initial compositions), the cooperative-basin size $\omega_{\mathrm C}$ (the
fraction of cooperative trajectories) and the cooperation level
$\bar\rho_{\mathrm C}$ within the basin, with
$\bar\rho\approx\omega_{\mathrm C}\bar\rho_{\mathrm C}$. For the total-depth
comparison, the envelope at total depth $K$ is the largest $\bar\rho$ among the
$K+1$ allocations with $k_1+k_2=K$.

\subsection*{Invasion of a discriminator population}

We consider a homogeneous DISC population at $r^*_{\mathrm R}$ and a rare
mutant $\mathrm X$ that uses the residents' norm and record mechanism. Because
the mutant is rare, it interacts with residents only and leaves their
reputations and assessments unchanged. Suppose that $\mathrm X$ intends to help
a resident recipient with probability $\chi_{\mathrm X}$, independently of the
observer's image of that recipient, and acts independently across interactions.
It is then assessed positively with probability
$u_{\mathrm X}=\chi_{\mathrm X}\Phi^*_{\mathrm C}+(1-\chi_{\mathrm X})\Phi^*_{\mathrm D}$,
its stationary reputation is $r^*_{\mathrm X}=\mathcal R(u_{\mathrm X})$, and its
normalized payoff advantage over residents is (Supplementary Note~7)
\begin{equation}
\Lambda_{\mathrm X}\equiv\frac{\Pi_{\mathrm X}-\Pi_{\mathrm R}}{c(1-e_1)}
=\bigl(r^*_{\mathrm R}-\chi_{\mathrm X}\bigr)
-\frac{b}{c}\bigl(r^*_{\mathrm R}-r^*_{\mathrm X}\bigr),
\label{eq:m-mutant-identity}
\end{equation}
the cooperation cost saved minus the benefit lost through reputation. The
resident itself corresponds to $\chi=r^*_{\mathrm R}$ and
$u=u^*_{\mathrm R}=r^*_{\mathrm R}\Phi^*_{\mathrm C}+(1-r^*_{\mathrm R})\Phi^*_{\mathrm D}$.

\textit{Unconditional mutants.} ALLD ($\chi=0$) has reputation
$\mathcal R(\Phi^*_{\mathrm D})$ and ALLC ($\chi=1$) has reputation
$\mathcal R(\Phi^*_{\mathrm C})$, where
$\Phi^*_{\mathrm D}=r^*_{\mathrm R}P^+_{\mathrm D\mid\mathrm G}+(1-r^*_{\mathrm R})P^+_{\mathrm D\mid\mathrm B}$
and $\Phi^*_{\mathrm C}$ analogously. When ALLD is regarded worse and ALLC
better than the residents, Eq.~\eqref{eq:m-mutant-identity} implies that the
discriminator population resists both mutants if and only if
\begin{equation}
R_{\mathrm L}\equiv
\frac{r^*_{\mathrm R}}
{r^*_{\mathrm R}-\mathcal R\bigl[r^*_{\mathrm R}P^+_{\mathrm D\mid\mathrm G}
+(1-r^*_{\mathrm R})P^+_{\mathrm D\mid\mathrm B}\bigr]}
<\frac bc<
\frac{1-r^*_{\mathrm R}}
{\mathcal R\bigl[r^*_{\mathrm R}P^+_{\mathrm C\mid\mathrm G}
+(1-r^*_{\mathrm R})P^+_{\mathrm C\mid\mathrm B}\bigr]-r^*_{\mathrm R}}
\equiv R_{\mathrm U}.
\label{eq:m-window}
\end{equation}
If ALLD is not regarded worse than the residents, no $b/c$ excludes it; if ALLC
is not regarded better, it cannot invade. Because
$u^*_{\mathrm R}=r^*_{\mathrm R}\Phi^*_{\mathrm C}+(1-r^*_{\mathrm R})\Phi^*_{\mathrm D}$,
the interval in Eq.~\eqref{eq:m-window} is non-empty if and only if
$r^*_{\mathrm R}\mathcal R(\Phi^*_{\mathrm C})+(1-r^*_{\mathrm R})\mathcal R(\Phi^*_{\mathrm D})
<\mathcal R(u^*_{\mathrm R})$, that is, if $\mathcal R$ at the resident's
assessment probability lies above the chord joining its values at
$\Phi^*_{\mathrm D}$ and $\Phi^*_{\mathrm C}$. Without inertia,
$\mathcal R$ is linear and $R_{\mathrm L}=R_{\mathrm U}=1/\Delta\Phi^*$, so no
benefit-to-cost ratio protects a discriminator population (Supplementary
Note~7). Numerically, we classify local outcomes from the signs of the
unnormalized payoff differences
$\lambda_{\mathrm{ALLD}}=\Pi_{\mathrm{ALLD}}-\Pi_{\mathrm R}$ and
$\lambda_{\mathrm{ALLC}}=\Pi_{\mathrm{ALLC}}-\Pi_{\mathrm R}$ (with $c=1$), which also
covers cases without finite thresholds: ALLD is not excluded if
$\lambda_{\mathrm{ALLD}}\ge-10^{-10}$; otherwise ALLC invades if
$\lambda_{\mathrm{ALLC}}>10^{-10}$; otherwise the residents resist both.

\textit{Opportunistic discriminators.} An opportunistic discriminator (O-DISC)
behaves like DISC, except that each intended cooperation towards a recipient it
regards as good is replaced, independently and before implementation error, by
intended defection with probability $\delta\in[0,1]$. Thus
$\chi=(1-\delta)r^*_{\mathrm R}$, and $\delta=0$ and $\delta=1$ recover DISC and
ALLD. The mutant's positive-assessment probability and reputation are
\begin{equation}
u_{\mathrm X}(\delta)=u^*_{\mathrm R}-\delta\,r^*_{\mathrm R}\,\Delta\Phi^*,
\qquad
r^*_{\mathrm X}(\delta)=\mathcal R\bigl[u_{\mathrm X}(\delta)\bigr],
\label{eq:m-odisc-reputation}
\end{equation}
and Eq.~\eqref{eq:m-mutant-identity} gives the normalized invasion payoff and
its slope at the resident,
\begin{equation}
\Lambda(\delta)=\delta\,r^*_{\mathrm R}
-\frac{b}{c}\bigl[r^*_{\mathrm R}-r^*_{\mathrm X}(\delta)\bigr],
\qquad
\Lambda'(0)=r^*_{\mathrm R}\Bigl[1-\frac{b}{c}\,\Delta\Phi^*\,
\mathcal R'(u^*_{\mathrm R})\Bigr].
\label{eq:m-odisc-invasion}
\end{equation}
A positive $\Lambda(\delta)$ means that the mutant can invade; $\Lambda(0)=0$,
and $\Lambda(1)$ is the normalized ALLD advantage. For $r^*_{\mathrm R}>0$,
$\Lambda(\delta)=\delta r^*_{\mathrm R}[1-(b/c)H(\delta)]$, where
$H(\delta)=[r^*_{\mathrm R}-r^*_{\mathrm X}(\delta)]/(\delta r^*_{\mathrm R})$ is the
reputation lost per unit of withheld cooperation, extended continuously by
$H(0)=\Delta\Phi^*\mathcal R'(u^*_{\mathrm R})$. The residents therefore exclude
every O-DISC mutant if and only if
\begin{equation}
\min_{0\le\delta\le1}\frac{r^*_{\mathrm R}-r^*_{\mathrm X}(\delta)}{\delta\,r^*_{\mathrm R}}
\;\ge\;\frac{c}{b},
\label{eq:m-security}
\end{equation}
with equality permitting neutral mutants. Because $H(\delta)$ equals
$\Delta\Phi^*$ times the slope of the chord of $\mathcal R$ between
$u_{\mathrm X}(\delta)$ and $u^*_{\mathrm R}$, for $\Delta\Phi^*>0$ the minimum
is attained at $\delta\to0$ whenever $\mathcal R$ is concave on that interval
(always for $k_2=0$), and Eq.~\eqref{eq:m-security} then reduces to
$\Lambda'(0)\le0$; in general, the whole interval must be checked. Near an
all-good resident with $k_1\ge1$, $\mathcal R'(u^*_{\mathrm R})\approx1-\alpha$,
so full tolerance readout makes weak opportunism cheap. In the error-free limit with $\beta=1$ and $k_1\ge1$, an
all-good resident excludes every O-DISC mutant if and only if
$\alpha\le1-c/b$, for every norm and every $k_2$ (Supplementary Note~8); at
$b/c=5$ this bound is the readout $\alpha=0.8$ used in Fig.~\ref{fig:opportunism}
and Extended Data Fig.~\ref{ed:opportunism-security}.

For fixed norm, depths, $\beta$ and errors, a readout $\alpha$ is secure if a
strictly stable high-reputation resident exists and satisfies
Eq.~\eqref{eq:m-security}. The lower boundary $\alpha_{\mathrm L}$ of a secure
interval, when it is not $\alpha=0$, is a fold of the high-reputation branch,
$F(r,\alpha)=\partial_rF(r,\alpha)=0$ with
$F(r,\alpha)=\mathcal R[u_{\mathrm R}(r);\alpha,\beta,k_1,k_2]-r$; the fold
itself is excluded and does not depend on $b/c$. The upper boundary
$\alpha_{\mathrm U}$, when it is not $\alpha=1$, is where
Eq.~\eqref{eq:m-security} holds with equality. In every case we computed it was
determined by $\delta\to0$, that is by
$(b/c)\,\Delta\Phi(r)\,\mathcal R'[u_{\mathrm R}(r);\alpha_{\mathrm U}]=1$
together with $F(r,\alpha_{\mathrm U})=0$. When the fold and the local condition
meet, the secure interval closes; this happens at
$b/c=u'_{\mathrm R}(r)/\Delta\Phi(r)$ evaluated at the fold, which equals 2 for
Shunning (Supplementary Note~8).

\textit{Alternative reputation filters.} For Fig.~\ref{fig:opportunism}a we
replace $\mathcal R$ by the stationary response of two alternative mechanisms
while keeping the norm, errors, resident selection and
Eqs.~\eqref{eq:m-odisc-reputation}--\eqref{eq:m-odisc-invasion} unchanged.
Aggregation, adapted from ref.~\cite{michelmata2024}, assigns a good image if at
least $Q$ of $M$ independently sampled assessments are positive,
$\mathcal R_{\mathrm{agg}}(u)=\sum_{j=Q}^{M}\binom Mj u^j(1-u)^{M-j}$, with
$(M,Q)=(2,1)$, so that $\mathcal R_{\mathrm{agg}}(u)=2u-u^2$. Quantitative
assessment, adapted from ref.~\cite{schmid2023}, keeps an integer score
$z\in\{-z_{\max},\dots,z_{\max}\}$ that moves up or down by one after a positive
or negative assessment (attempts to leave the range leave the score unchanged)
and is read as good if $z\ge z_0$; with $t=u/(1-u)$,
$\mathcal R_{\mathrm{score}}(u)=\sum_{j=z_{\max}+z_0}^{2z_{\max}}t^j\big/\sum_{j=0}^{2z_{\max}}t^j$,
and we use $z_{\max}=5$, $z_0=0$. Neither mechanism was optimized or matched in
memory to the inertial record.

\subsection*{Finite-population reputation simulations}

Individual-based simulations test reputation formation without the mean-field
closure, at fixed strategy composition; they do not simulate strategy
evolution. A population of $N$ individuals contains fixed numbers $n_I$ of each
strategy. Every DISC individual stores a latent record of every individual,
including itself, and all records start in $\mathrm G$. Records held by ALLC and
ALLD individuals are not stored because they never affect actions or
assessments. In each event: (i) a donor is drawn uniformly from the population
and a recipient uniformly from the remaining $N-1$ individuals, and an observer
is drawn independently and uniformly from the DISC individuals (the observer
may coincide with donor or recipient); (ii) the donor forms its intended
action, a DISC donor by sampling its own image of the recipient from
Eq.~\eqref{eq:m-readout}; (iii) implementation error is applied; (iv) the
observer independently samples its own image of the recipient, applies the norm
to the realized action and reverses the assessment with probability $e_2$; and
(v) the observer's record of the donor is updated. A sweep consists of
$n_{\mathrm{DISC}}N$ events, one update per stored record on average.

At sampled times we use $s_{ij}=\sigma(L_{ij})$, which averages exactly over
readout randomness. For target $j$, the mean goodness and the probability that
two distinct DISC observers give different images are
\begin{equation}
g_j=\frac{1}{n_{\mathrm{DISC}}}\sum_i s_{ij},\qquad
d_j=\binom{n_{\mathrm{DISC}}}{2}^{-1}\sum_{i<i'}
\bigl[s_{ij}(1-s_{i'j})+(1-s_{ij})s_{i'j}\bigr],
\label{eq:m-finite-estimators}
\end{equation}
where $i,i'$ run over DISC observers. Averaging over type-$I$ targets gives
$g_I$ and $d_I$, which are then averaged over sampling times. The reputational
separation and mean disagreement are
\begin{equation}
\Delta g=\frac{g_{\mathrm{DISC}}+g_{\mathrm{ALLC}}}{2}-g_{\mathrm{ALLD}},
\qquad
\bar d=\sum_I\frac{n_I}{N}\,d_I .
\label{eq:m-recognition}
\end{equation}
Unlike the mean-field theory, the simulations retain target-specific variation
and correlations between the private records of different observers.

The recognition analyses use equal strategy numbers, $\alpha=\beta=1$ and
$0\le k_1,k_2\le6$. Each run lasts 1,000 sweeps and is sampled every five sweeps
from sweep 800 to sweep 1,000 (41 samples). The depth maps in
Fig.~\ref{fig:recognition}e--h use one run per condition with $N=150$. The
replicated analysis (Extended Data Fig.~\ref{ed:recognition-selection}) uses 30
independent runs per condition with $N=150$ and reports means across runs, with
pointwise 95\% percentile intervals from 2,000 bootstrap resamples of the run
means. A condition is eligible for selection if, for both $\Delta g$ and
$\bar d$, the bootstrap interval of the difference between the second- and
first-half means of the sampling window lies within $(-0.01,0.01)$. Two slowly
relaxing SJ conditions required longer runs: $(k_1,k_2)=(4,5)$ with 3,000 sweeps
(sampled from sweep 2,400) and $(3,5)$ with 50,000 sweeps (sampled from sweep
40,000); all 196 conditions then passed. Among eligible conditions with
$\Delta g\ge0.8$ and $\bar d\le0.1$, we select the smallest $K$, breaking ties
by larger $\Delta g$, then smaller $\bar d$, $k_1$ and $k_2$. This gives the
four architectures listed above. The stars in Fig.~\ref{fig:recognition}e--h
apply the same cutoffs and tie-breaking rule to the single-run maps, without
the drift criterion, and identify the same four architectures. Selection uses
reputation statistics only.
The robustness analysis varies the separation cutoff from 0.70 to 0.90 and the
disagreement cutoff from 0.050 to 0.150 in steps of 0.05 and 0.025. The
illustrative image matrices in Fig.~\ref{fig:recognition}a--d are final states of
single runs with $N=90$ after 1,000 sweeps.

\subsection*{Numerical procedures and figure-specific settings}

All panels use $b/c=5$ and $e_1=e_2=0.02$ unless an axis varies them;
Supplementary Note~12 (Supplementary Table~S11) lists every setting, grid and
tolerance.
$\mathcal R$ and $\mathcal R'$ are evaluated in closed form. Homogeneous-DISC
roots (Fig.~\ref{fig:opportunism}; Extended Data
Figs.~\ref{ed:resident-reputation}a--d, \ref{ed:reputation-planes} and
\ref{ed:opportunism-security}) are located by evaluating $F(r)$ on a fine
uniform grid, refining sign changes by bracketing and classifying roots by the
sign of $F'(r)$. At mixed compositions, Eq.~\eqref{eq:m-relaxation} is
integrated from $r_I=1$ by fourth-order Runge--Kutta steps until
$\max_I|\mathcal R(u_I)-r_I|\le10^{-11}$ ($10^{-12}$ for
Fig.~\ref{fig:strategy-dynamics}). For local invasion (Extended Data
Figs.~\ref{ed:benefit-cost}--\ref{ed:readout-evolution}), the resident
relaxation is followed by step-limited Newton refinement, which keeps the
branch reached from $r=1$, to residual $\le10^{-13}$.

All sampled-cooperation results (Fig.~\ref{fig:memory-allocation}; Extended
Data Figs.~\ref{ed:resident-reputation}e, \ref{ed:benefit-cost},
\ref{ed:depth-evolution} and \ref{ed:readout-evolution}e--h) use one engine.
Selected reputations are tabulated exactly on a triangular lattice with 120
subdivisions and interpolated linearly only inside cells in which no reputation
component varies by more than 0.02; elsewhere, including every cell crossed by
a fold, they are computed exactly at the requested composition. Payoffs and the
replicator field are always computed from these reputations, never
interpolated. Trajectories are integrated with the explicit midpoint rule
($\Delta t=0.2$) until $\rho_{\mathrm{int}}$ changes by at most $10^{-7}$ over
each of three consecutive intervals of 100 time units, or for at most
$4\times10^4$ time units. Halving the step, refining the lattice and tightening
the tolerances changed $\bar\rho$ by at most $1.3\times10^{-5}$ and no outcome
class. Initial compositions are 512 points of a randomly shifted Halton
sequence (bases 2 and 3; seed 2026) or, for Extended Data
Fig.~\ref{ed:benefit-cost}, eight independently scrambled 512-point Sobol
sequences~\cite{virtanen2020} (seeds 2026--2033), mapped area-uniformly onto the
simplex by reflecting points of the unit square with $x+y>1$. In Extended Data
Fig.~\ref{ed:benefit-cost}, error bars are standard errors across the eight
samples and therefore describe the sampling of initial compositions; the
dynamics are deterministic. Extended Data Fig.~\ref{ed:resident-reputation}e
initializes reputations once, at $t=0$, as all-good or all-bad and then follows
them quasi-statically along each trajectory by warm-started relaxation, so that
a branch is left only where it ends.

In Fig.~\ref{fig:strategy-dynamics} and Extended Data
Fig.~\ref{ed:readout-evolution}a--d, reputations are relaxed exactly at every
point of an $84\times74$ grid clipped to the simplex. Vertex stability follows
from the two rare-strategy growth rates; edge equilibria are located on
501-point edge scans, and their markers refer to stability along the edge;
interior equilibria are classified by their Jacobian. Under SC without inertia,
all three payoffs are equal on the line
$f_{\mathrm{DISC}}=c/[b(\epsilon-e_2)]\approx0.2126$ (Supplementary Note~6).
Extended Data Fig.~\ref{ed:reputation-planes} reports the highest stable
homogeneous-DISC root and marks bistability (one-step readout plane,
$161\times161$; depth plane, $0\le k_1,k_2\le8$); for every root in the depth
planes, stability in Eq.~\eqref{eq:m-relaxation} agreed with the linearization of
the full latent-state dynamics. On the ALLD--DISC edge (Extended Data
Fig.~\ref{ed:edge-geometry}), Eq.~\eqref{eq:m-self-consistency} reduces exactly
to one equation in $\bar r$, which is solved for 1,071 values of
$f_{\mathrm{DISC}}\in[0.01,1]$. Local invasion maps use $241\times241$ points of
$e_1,e_2\in[0,0.45]$ (Extended Data Fig.~\ref{ed:error-invasion}) and
$801\times801$ points of $(\alpha,\beta)$ (Extended Data
Fig.~\ref{ed:readout-evolution}i--l). In Fig.~\ref{fig:opportunism},
$\Lambda(\delta)$ is evaluated at 5,015 values of $\delta$, concentrated near
$\delta=0$, with sign changes refined by Brent's method; $\delta_{\mathrm c}$ is
the upper end of the interval with $\Lambda>0$. Residents of the alternative
filters are obtained from $r=1$ by damped iteration,
$r\leftarrow0.8r+0.2\,\mathcal R_{\mathrm X}[u_{\mathrm R}(r)]$ with
$\mathcal R_{\mathrm X}\in\{\mathcal R_{\mathrm{agg}},\mathcal R_{\mathrm{score}}\}$.
For Extended Data Fig.~\ref{ed:opportunism-security} ($\ge601$ values of
$\alpha$, 161 values of $b/c\in[2,8]$, $\beta=1$), $\min_\delta H(\delta)$ is
found by comparing $H$ at $\delta=0$, $\delta=1$ and all real roots in
$(10^{-8},1)$ of the numerator of the rational function
$\mathrm dH/\mathrm d\delta$, and a readout is classified as secure if
$1-(b/c)\min_\delta H\le10^{-10}$. Fold endpoints are excluded; endpoints at the
domain boundary or at neutral invasion are included.

\subsection*{Data and code availability}

The code, numerical data and figure scripts used in this study are provided
with the paper (verified computational package; one notebook per figure, a C
engine for the strategy dynamics, input configurations, validation records and
checksums). \textit{[Repository and DOI to be added before submission.]}
